\documentclass[12pt]{report}
\usepackage[portuguese]{babel}
\usepackage[utf8]{inputenc}
%\usepackage[dvips]{graphicx}
\usepackage{graphicx}
\usepackage{enumitem}
\usepackage[normalem]{ulem}
\usepackage{latexsym}
\usepackage{amsmath}
\usepackage{float}
\usepackage{xurl}
\usepackage{afterpage} % para páginas em branco
\usepackage{units} % Para frações bonitas
\usepackage{verbatim} % Para que o bloco {comment} funcione
\usepackage[pdfborder={0 0 0}]{hyperref}
\usepackage{wrapfig} % Para o qrcode do editorial
\usepackage[labelformat=empty]{caption} % Para tirar o "Figure:" da caption do QRCode

%\usepackage{showframe} %Para mostrar o frame do arquivo

\newcounter{qcounter}

\setlength{\parindent}{0pt}
\setlength{\parskip}{3ex plus 0.5ex minus 0.5ex}

\addtolength{\voffset}{-2.8cm}
\addtolength{\textheight}{5.5cm}
\addtolength{\hoffset}{-2.3cm}
\addtolength{\textwidth}{4.5cm}
\setlength{\topmargin}{0cm}

%FIXME GAMBIARRA
%\setlength{\evensidemargin}{55pt}

\pagestyle{plain}

% =============== Seção de definições de macros ========================

% Delimita uma seção:
\newenvironment{secao}[1] {
    \framebox[\textwidth] {
        \rule[-1.2ex]{5ex}{5.5ex}
        {\Large\sf #1}
        \hspace{\stretch{1}}
    }
    \phantomsection %Faz o link do pdf funcionar direito
	\addcontentsline{toc}{chapter}{#1}
    \nopagebreak[4]
}{}

%FIXME Isso pode ser resolvido de um jeito melhor, fazendo uma seção (editorial e mandamentos)
% não entrarem no índice usando uma tag, e não outro tipo de environment.
% Veja a primeira linha de editorial.tex para entender melhor.
% Delimita o editorial (que não entra no Índice):
\newenvironment{editorial}[1] {
    \framebox[\textwidth] {
        \rule[-1.2ex]{5ex}{5.5ex}
        {\Large\sf #1}
        \hspace{\stretch{1}}
    }
    \phantomsection %Faz o link do pdf funcionar direito
    %\addcontentsline{toc}{chapter}{#1}
    \nopagebreak[4]
}{
\thispagestyle{empty}
\pagebreak
}


% Delimita uma subseção
\newenvironment{subsecao}[1] {
    \rule[0ex]{2.5ex}{2.5ex}
    {\Large\sf #1}
    \phantomsection %Faz o link do pdf funcionar direito
	\addcontentsline{toc}{section}{#1}
    \nopagebreak[4]
}{}

% Delimita uma subsubseção
\newenvironment{subsubsecao}[1] {
    \rule[0ex]{1ex}{1.5ex}
    {\large\sf #1}
    \phantomsection %Faz o link do pdf funcionar direito
% Próxima linha é responsável por subsubseções não aparecerem no índice!
	%\addcontentsline{toc}{subsection}{#1}
    \nopagebreak[4]
}{}

% Coloca uma figura grande (sem ser quadrinhos)
\newcommand{\figuragrande}[1] {
    \begin{figure}[!htbp]
      \begin{center}
        \includegraphics[width=\textwidth]{img/#1.pdf}
      \end{center}
    \end{figure}
}

% Coloca uma figura menor (sem ser quadrinhos)
\newcommand{\figurapequenainline}[1] {
  \begin{wrapfigure}{r}{0.25\textwidth}
    \vspace{-25pt}
    \begin{center}
      \includegraphics[width=0.25\textwidth]{img/#1.pdf}
    \end{center}
    \vspace{-25pt}
  \end{wrapfigure}
}

% Coloca uma figura menor (sem ser quadrinhos) apertada verticalmente
\newcommand{\figurapequenainlineflexivel}[2] {
  \begin{wrapfigure}{r}{0.25\textwidth}
    \vspace{-#2}
    \begin{center}
      \includegraphics[width=0.25\textwidth]{img/#1.pdf}
    \end{center}
    \vspace{-#2}
  \end{wrapfigure}
}

% Coloca uma figura menor (sem ser quadrinhos) apertada verticalmente
\newcommand{\figurapequenainlineapertada}[1] {
  \begin{wrapfigure}{r}{0.25\textwidth}
    \vspace{-40pt}
    \begin{center}
      \includegraphics[width=0.25\textwidth]{img/#1.pdf}
    \end{center}
    \vspace{-40pt}
  \end{wrapfigure}
}

% Coloca quadrinhos
\newcommand{\quadrinhos}[1] {
    \figuragrande{quad#1}
}

% Coloca um mapa
\newcommand{\mapa}[1] {
	\begin{figure}[H]
		\centering
		\includegraphics[width=\textheight, angle=90]{img/#1}
	\end{figure}
}

%FIXME GAMBIARRA
% Coloca um mapa do IME virando ele do jeito certo
\newcommand{\mapaime}[1] {
    \begin{figure}[H]
        \centering
        \includegraphics[width=\textwidth, angle=180]{img/#1.pdf}
    \end{figure}
}

% Comando auxiliar para o \paginaembranco
\newcommand\blankpage{%
    \null
    \thispagestyle{empty}%
    \addtocounter{page}{-1}%
    \newpage
    }

% Cria uma página em branco SEM NUMERAÇÃO após completar a atual
% (assim que é possível)
\newcommand\paginaembranco{%
    \afterpage{\blankpage}
    }

% ============================ Documento ===============================
\begin{document}

% Capa -------------------------------------------------------------------------
\begin{figure}[p]
  %\begin{center}
    %\hspace{1.7cm}
    \includegraphics[height=1\textheight]{img/capa2025.pdf} %REFTIME
  %\end{center}
\end{figure}
\thispagestyle{empty} % Some com aquele número 1 feio na capa.
\clearpage
\newpage

% Começa a numeração aqui para que o LaTeX entenda qual
% página é ímpar e qual é par antes do índice.
\pagenumbering{arabic}

% Pré-Índice -------------------------------------------------------------------
\begin{editorial}{Editorial}

Bixes, foi difícil chegar até aqui. Vocês
estão meio ou completamente perdidos. Temos apenas uma sugestão: aproveitem esta
etapa! Façam da sua estadia na USP a melhor época das suas vidas! Vocês verão
que a USP tem muitas e muitas coisas a oferecer! Não se preocupem apenas em
estudar e passar de ano, como vocês fizeram durante sua vida inteira; aproveitem
TUDO! Vocês podem não acreditar nisso agora, mas saibam que viverão momentos
inesquecíveis aqui no IME: alguns fantásticos, outros deploráveis.

Este guia foi feito para que vocês (bixes que não gostam de ler, apenas olhem as figuras)
possam aprender um pouquinho do que é a USP,
o IME e a vida universitária que se inicia agora. 
Gostaríamos que guardassem este guia com muito carinho para que, futuramente, possam consultá-lo
quando surgir alguma dúvida do que fazer em determinada situação!

O guia foi escrito numa forma
descontraída e fácil para que vocês consigam entender; mesmo assim, se aparecer
alguma dúvida, vocês podem se dirigir a qualquer veterane, e sua dúvida será
sanada (e, quem sabe, talvez você também comece uma nova amizade).

Além disso, ao invés de usarmos as palavras ``bixos'' ou ``bixetes'', vamos nos referir
a vocês como ``bixes'', pois estamos dirigindo palavra a todas as identidades de 
gênero! Por motivos de vocabulário e acessibilidade, usamos a linguagem neutra 
apenas em ``bixes'' e ``veteranes''.

%Empurra pro fim da página (altura)
%\vspace{\stretch{1}}

%Linha divisória
\rule{\textwidth}{0.5ex}\rule{2ex}{0.5ex}

%REFTIME
{\large\bf Guia de Bixe 2025} \\
Uma publicação da Comissão de Recepção


\begin{wrapfigure}{R}{0.25\textwidth}

% Link para "https://ime.usp.br/~kazuo/guia"
% Esse link contém um contador de acessos que redireciona para o guia correto
% Em 2020: https://ime.usp.br/~kazuo/guiadobixo.pdf
  \begin{center}
    \includegraphics[width=0.22\textwidth]{img/qrcodeguia2024.png}
  \end{center}
  \vspace{-20pt}
  \caption{Baixe aqui o Guia em formato digital!}
  \vspace{80pt}
\end{wrapfigure} 

Muitas pessoas dedicaram seu tempo (e suas férias) para que esse Guia ficasse
pronto: desde o Donald Knuth em 1978, na criação do \TeX\makebox{}, até o
pessoal da Comissão na correria de ontem à noite. Esperamos que vocês gostem do Guia
e adotem-no como livro de cabeceira. Por fim, nos avisem de qualquer informação
incorreta ou desatualizada, afinal vocês também são responsáveis por tudo que o
IME oferece a partir de agora.

Lembrem-se: este é o primeiro e último ano de vocês como bixes. Aproveitem!
%\vspace{2em}

%REFTIME (Capa muda todo ano)
%Capa de 2022: ???
%Capa de 2021: ???
%Capa de 2020: ???
%Capa de 2019: ???
%Capa de 2018: ???
%Capa de 2017: ???
%Capa de 2016: public domain
%Capa de 2015:
%{\bf Imagem da Capa:} Viktoria Ridzel -- \href{www.viria13.deviantart.com}{viria13.deviantart.com}

Este trabalho está licenciado sob a Licença Attribution-ShareAlike 4.0 da
Creative Commons. Para ver uma cópia desta licença, visite
\url{https://creativecommons.org/licenses/by-sa/4.0/deed.pt_BR} ou envie
uma carta para Creative Commons, 444 Castro Street, Suite 900, Mountain View,
California, 94041, USA.
\\
\begin{figure}[H]
    \centering
    \includegraphics{img/cc/by-sa.png}
\end{figure}

\end{editorial}


% Carta Aos Ingressantes -------------------------------------------------------
\input{carta_aos_ingressantes.tex}

% Os Dez Mandamentos -----------------------------------------------------------
\begin{editorial}{Os Dez Mandamentos}

\begin{enumerate}
  \item Os veteranes têm sempre razão, exceto quando os bixes têm;
  \item Na hipótese de o bixe ter razão, entra imediatamente em vigor
        o primeiro mandamento;
  \item Em qualquer evento social, as despesas correm sempre por conta dos
        bixes, menos o que os veteranes consomem;
  \item Os bixes têm o direito de permanecerem calados (exceto quando
        interpelados por veteranes ou quando quiserem falar). Tudo o que eles
        disserem pode ser usado contra eles, \textit{c’est la vie};
  \item Os bixes devem se apresentar imediatamente em caso de convocação por
        veteranes, a não ser que estejam fazendo algo melhor da vida. Os bixes
        que não se apresentarem se arrependerão amargamente por perderem
        conselhos valiosos; %não queremos ameaçar ninguém :)
  \item Não são garantidos no IME os direitos constitucionais dos bixes à vida
        social, liberdade de sono e igualdade de notas;
  \item Os bixes devem estar prontos para assumirem as seguintes funções para
        veteranes: completar mesa de jogo, companhia de bandejão, completar
        time(s) etc., mas somente quando tiverem vontade;
  \item Os bixes devem amar e respeitar seus veteranes, a vida, as plantas,
        todo mundo e acima de qualquer coisa, Goku;
  \item Para os casos não abrangidos por estas regras, a decisão final correrá
        por conta dos veteranes, com exceção do que não lhes diz respeito;
  \item Todos os bixes estão despreparados, exceto os que leram isto.
\end{enumerate}

Como bixes, vocês têm todo o direito de reclamar dos mandamentos! 
Qualquer reclamação deverá ser protocolada em três vias datadas, assinadas e
autenticadas, com firma reconhecida em cartório, e assim encaminhadas à
Comissão de Recepção 2025 %REFTIME
via mala direta e serão imediatamente incineradas. Alguns veteranes sugeriram
que incinerássemos os reclamantes também. Ao invés disso, cogitamos incinerar
tais veteranes. A medida passou por estudo e, devido aos custos operacionais e
ao lixo tóxico que seria produzido, foi rejeitada.

OU vocês podem conversar diretamente com um veterane da Comissão :)

\end{editorial}


% Os Sete Pecados --------------------------------------------------------------
\input{sete_pecados.tex}


%Renomeando o índice------------------------------------------------------------
\renewcommand{\contentsname}{\center Esse guia contém...}

% Coloca o índice --------------------------------------------------------------
\tableofcontents
% Remove números de página do índice
\thispagestyle{empty}
\addtocontents{toc}{\protect\thispagestyle{empty}}
\addtocontents{toc}{\protect\pagestyle{empty}}
\newpage
\thispagestyle{empty}
\cleardoublepage

% Começa a numeração aqui
\pagenumbering{arabic}


% Comissão de Trote e o Kit-bixe -----------------------------------------------
\begin{secao}{Comissão de Recepção e \textit{Kit}-bixe}

A Comissão de Recepção de Calouros, também conhecida apenas como Comissão de
Recepção, é responsável por auxiliar es ingressantes em seus primeiros momentos
IMEanos. Sabemos que não é um momento fácil. Vocês estão entrando em uma nova fase
de suas vidas, em um lugar estranho, com pessoas estranhas (em todos os sentidos),
e nosso objetivo é fazer com que vocês se sintam bem-vindes e se integrem
($\int$) com coleguinhas e veteranes!

A Comissão de Recepção organiza a super Semana de Recepção, cheia de atividades
legais! (Vejam a programação em: \url{https://recepcao.ime.usp.br/}).
Somos nós também que organizamos o magnífico encontro des bixes: AnIME-se.
Atente-se à página da Comissão no Facebook, TikTok ou Instagram ou em um dos
grupos de WhatsApp (BIXES 2025 IME USP) para quando a data for divulgada. %REFTIME

A Comissão de Recepção é formada peles mais animades e divertides veteranes. Como
dissemos, organizam a Semana de Recepção, ajudam bixes perdides, mantêm um clima
alegre durante a recepção, fazem isso e aquilo...
Vocês devem estar pensando ``Puxa! Como nós, bixes, podemos retribuir
tamanha dedicação?'' É simples, bixes: {\bf\em comprem o \textit{kit}-bixe}!!!

O \textit{kit}-bixe, como vocês devem saber, é um conjunto de coisas
importantíssimas para vocês, ingressantes perdides! Ele contém dois tipos de
itens:
\begin{itemize}
\item itens úteis;
\item itens essenciais.
\end{itemize} %REFTIME (Isso muda todo ano)
Dentre eles, temos uma camiseta com uma linda estampa do tema do ano de vocês,
que serve para que nós os identifiquemos como bixes e para que o pessoal
ache que vocês são inteligentes; um lindo estojo com o símbolo do IME e vários 
materiais para vocês usarem nas aulas e mostrarem por aí; uma caneca e 
um talabarte MARAVILHOSOS (também personalizados) para vocês economizarem muitos 
copos no bandejão; um chaveiro personalizado muito maneiro que também serve para 
abrir garrafas; dois adesivos, um para vocês colarem no carro que vão ganhar de 
presente, e outro para os papais (quando vocês virem o adesivo vocês vão entender); 
um caderno super lindo e sem pauta para vocês desenharem seus gráficos das aulas 
de Cálculo; uma tatoo de henna e um cordão de cartão de metrô. 
Tudo isso preparado com muito carinho e especialmente para vocês pela 
Comissão de Recepção, além de estarem contidos em uma mochila sport do IME-USP, 
afinal, não dá pra carregar tudo isso na mão, né?

Além dessa maravilha de \textit{kit}-bixe, vocês ainda podem adicionar à parte bottons
exclusivos e personalizados com o símbolo do IME-USP e com o tema da recepção.

Adquiram o maravilhoso \textit{kit}-bixe do IME-USP! Fiquem atentes às redes sociais da
Comissão de Recepção para saberem quais os períodos em que as vendas estarão abertas.
O kit estará à venda na Semana de Recepção, mas somente enquanto
durarem os estoques. (Não vão chorar depois, hein?)

Para terminar, além de tudo isso, a Comissão de Recepção é que faz esse maravilhoso 
guia que vocês estão lendo agora (ou estão só olhando as figuras, vai saber...). 
Esperamos que vocês gostem do nosso trabalho! Qualquer coisa, nos
procurem! Entrem na nossa página no Facebook (\url{https://www.facebook.com/recepcaoimeusp}),
no nosso TikTok (\url{https://tiktok.com/@imeusp.recepcao})
ou no nosso Instagram (\url{https://instagram.com/imeusp.recepcao/}) contando
o que vocês sentiram ao ler o guia, o que fizeram de errado com vocês na
recepção (com o nome do veterane na denúncia), o que vocês acharam da Semana
de Recepção e se vocês se sentiram bem-vindes, talvez seja feito um forms para 
vocês nos avaliarem. Estaremos sempre prontes a ajudá-les!
\end{secao}

\quadrinhos{1}
\pagebreak

% O CAMat ----------------------------------------------------------------------
\begin{secao}{O CAMat}

CAMat é a sigla para Centro Acadêmico da Matemática, Estatística e Computação 
"Elza Furtado Gomide". Bom, o que faz um centro acadêmico? 
O centro acadêmico é a entidade representativa des estudantes (conhecido 
também como entidade de base), ele serve para organizar o movimento estudantil
dentro do Instituto, lutando por melhorias internas no IME, e também
participando de assuntos da política da USP e de fora dela.

No caso do CAMat, representamos todes aquelus matriculades na graduação e pós-graduação do IME-USP.

O administrativo da entidade fica no interior da Vivência (sala 18 do bloco B), lá você 
poderá encontrar uma vendinha de guloseimas e bebidas para matar a larica e sede nossa de
cada dia. Como a Vivência,  o acesso é livre e aberto a todos! Além dos
comes e bebes, também emprestamos - mediante a apresentação do cartão USP -
calculadoras científicas, baralhos e outros jogos de carta! Na Vivência em si, temos sofás
para vocês passarem um tempo, tirar um cochilo, jogar sinuca ou Smash, e
também armários para alugar. Enfim, a Vivência + CAMat são o espaço de convivência 
des estudantes, com a manutenção e organização de responsabilidade do centro acadêmico 
- o que não exclui a necessidade de qualquer pessoa que frequente o espaço também se 
empenhar na conservação da limpeza e ordem na Vivência.

Importante reforçar que a atuação do CAMat não se basta à salinha ou à Vivência. 
O centro acadêmico é responsável por diversos projetos - a exemplo, CinIME e Aulões 
-, e também promove eventos para maior integração entre os estudantes e o espaço do 
instituto e da universidade. Aliás, caso você tenha propostas de eventos ou projetos, 
fale conosco, dado que sendo o meio des estudantes agirem e se organizarem na luta pela 
universidade popular, estamos sempre abertos ao debate de ideias e a fomentação de discussões 
construtivas. O diálogo com es estudantes é essencial para que a gestão seja ativa e presente no IME!

Veja a seguir um pouco sobre os projetos do CAMat:


\begin{subsecao}{Banco de Provas}

O CAMat, conjuntamente com es estudantes, organizam um Banco de Provas
\href{https://camat.ime.usp.br/apoio/} com VÁAAARIAS provas dos anos 
anteriores feitas pelos maravilhosos veteranes.

Não se esqueça de contribuir com sua prova ao final do semestre, não
deixe o Banco de Provas morrer! 

\end{subsecao}

%\begin{subsecao}{Aulões do CAMat e Oficina de Demonstrações}

%Os aulões são atividades promovidas pelo CAMat para ajudar os bixes com as
%disciplinas do primeiro ano do curso. Muitas vezes algum veterane que tem mais
%experiência em certa disciplina se oferece para dar uma aula tirando as dúvidas
%mais comuns e tentando passar um ponto de vista diferente para complementar e facilitar
%as aulas. As vezes os professores do IME podem ajudar também, como é o caso da Oficina 
%de Demonstrações que foi ministrada por uma professora para ajudar os bixes a entenderem 
%melhor como se faz uma prova matemática rigorosa, e até se tornou uma bolsa PUB para que alunos
%possam estudar mais sobre o uso dessa ferramenta pedagógica. 
%Participem dos aulões e usem esse espaço para tirar suas dúvidas e complementar as aulas!

%\end{subsecao}

\begin{subsecao}{CinIME}

Com a proposta de pautar de maneira mais crítica quais, como e onde consumimos 
entretenimento audiovisual, o CAMat organiza as atividades de cinema no IME 
como parte de um momento de reflexão, mas também de lazer. Afinal, cinema e 
pipoquinha são legais e todo mundo gosta! Toda sexta-feira são exibidos filmes 
diferentes com pipoca e refrigerante de graça, e as sessões são abertas a todes 
es alunos, funcionários e docentes. O CinIME também tem um espaço exclusivo nesse
guia para que vocês possam saber mais detalhes.

\end{subsecao}

\begin{subsecao}{Chá Mate e BoletIME}

Chá Mate é o encontro mensal que o CAMat organiza para es estudantos debaterem temas
ligados a política e a vida universitária. Desde a criação do Chá Mate, já discutiu-se
desde a falta de água no CRUSP até a elitização do acesso ao futebol (esse tema com parceria
da Atlética). Sempre buscamos trazer um material de apoio, que pode ser um texto, podcast, matéria
de jornal ou qualquer outra coisa que seja útil para abordar o assunto. 
Além disso, a organização do Chá Mate edita o BoletIME, uma espécie de jornal des alunes do IME em
que publicamos entrevistas, notícias sobre a USP, editoriais políticos, repasses dos Representantes
Discentes e novidades do Centro Acadêmico. Toda a comunidade estudantil IMEana é bem vinda a construir
o BoletIME conosco: existe um forms para enviar textos, e você pode falar sobre qualquer assunto da sua
experiência universitária. Você pode conferir a última edição do BoletIME em formato 
digital no site do CAMat, e em formato impresso na salinha do CAMat na vivência.

\end{subsecao}


\begin{subsecao}{An(IME)\texorpdfstring{$^2$}{²}}

An(IME)$^2$ é o evento de animes e mangás do IME, organizado pelo CAMat. Entre as atrações,
contamos com PokéBingo, Just Dance, Karaokê, concurso de Cosplay e de Cospobre e uma sessão 
especial do CinIME.
Também é quando você pode comprar o já tradicional Kit Otaku, que normalmente vem com uma
caneca, uma camiseta e um caderno temáticos de anime, além dos bottons exclusivos do evento.
É um dia para ê IMEane liberar o seu lado otaku e correr solto pelo Instituto depois de
encher a cara de Mupy!

\end{subsecao}

\begin{subsecao}{Campeonato de Sinuca}

Se você gosta de jogar sinuca, cola na Vivência que a mesa está sempre disponível
e basta pedir para jogar para ter sua vez nela. Se tem interesse em aprender a jogar,
não se sinta intimidade pelas pessoas que já estão na mesa, elas já passaram pelo
que você está passando e estão dispostas a te ajudar a aprender.

Fora as partidinhas entre aulas, organizamos também um campeonato de sinuca. Costuma
ter uma taxa de inscrição, parte do dinheiro é revertida em prêmio e a outra parte
destinada à manutenção da mesa.

Em 2023, a mesa passou por uma reforma que a deixou novinha em folha, então aproveitem
ao máximo esse momento para socializar e descansar um pouco no espaço do CAMat! %REFTIME

\end{subsecao}


\begin{subsecao}{Aluguel de Armário}

Agora, para salvar suas costas dos milhares de livros de Defesa Contra as Artes
das Trevas (Cálculo, Álgebra, Análise...) que você precisará carregar, o CAMat
também aluga armários mediante uma taxa relativamente insignificante (fiquem de olho
nas chamadas pra alugar armários). A distribuição costuma ocorrer através de sorteio -
podendo escolher entre armário alto ou baixo -, de forma que é necessário ficar atente
às datas do período de inscrição e chamada do sorteio! Confira seu e-mail para ficar por 
dentro dos prazos.

\end{subsecao}

\begin{subsecao}{Projeto Conectividade}

Durante a pandemia (2020 e 2021), dado a adoção do ensino remoto, o centro acadêmico em conjunto 
com a comunidade IMEana criou o Projeto Conectividade para mapear es estudantes que não tinham acesso
às aulas por falta de equipamento adequado ou conexão. Com esse mapeamento, conseguimos distribuir 
computadores para aquelus que precisavam. Em relação a conexão, trabalhamos em conjunto com a diretoria, 
que organiza a distribuição dos kits de internet disponibilizados pela Reitoria.

\end{subsecao}

O CAMat também tem várias formas de você entrar em contato sempre que precisar:

\begin{description}
\item [E-mail:] camat@ime.usp.br
\item [Site:] \url{https://www.ime.usp.br/~camat/}
\item [Telegram:] \url{https://t.me/camat_usp}
\item [Instagram:] \url{https://instagram.com/camat.usp}
\item [Whatsapp:] \url{https://chat.whatsapp.com/DZvJKoUXZ8247LmerFsRUS}
\end{description}

\textbf{LEMBREM-SE:} OCUPE PRA CARALHO a salinha do CAMat, a Vivência e a USP inteirinha, esses são espaços seus também!


\end{secao}

\quadrinhos{2}
\pagebreak

% A Atlética -------------------------------------------------------------------
\begin{secao}{A Atlética} %REFTIME - Quase tudo aqui é REFTIME...

\begin{subsecao}{O que é a AAAMat?}

É a Associação Atlética Acadêmica da Matemática - a entidade mais divertida do
mundo!!!! - que tem como objetivo trazer os melhores momentos da sua vida
universitária! A Atlética é formada por um grupo de IMEanes (gestão) que é
responsável por organizar atividades esportivas e eventos (festas, pizzadas,
premiações etc) para a comunidade IMEana.

\end{subsecao}
%%%%%%%%%%%%%%%%%%%%%%%%%%%%%%%%%%%%%%%%%%%%%%%%%%%%%%%%%%%%%%%%%%%%%%%%%%%%%%%%

Algumas das atividades da Atlética:

\begin{subsecao}{Atividades esportivas internas:}

Na AAAMat existem diretores de modalidade (DMs) que são pessoas responsáveis
pelos treinos e campeonatos dos seguintes esportes: futebol de campo, futsal,
basquete, vôlei, handebol, atletismo, natação, tênis de mesa, tênis de campo,
xadrez, baseball, softball, bridge, sinuca, truco, parkour, judô, karatê,
kendô, jiu-jitsu, rugby e e-sports (\textit{League of Legends}, \textit{Osu!},
\textit{Counter Strike: Global Offense}, \textit{Brawl Stars}, \textit{FIFA}
\textit{Rocket League}, \textit{Clash Royale}, \textit{Valorant},
\textit{Call of Duty: mobile}, \textit{Just Dance}, \textit{Hearthstone},
\textit{Pokemon UNITE}, \textit{Rainbow 6}, \textit{Super Smash Bros},
\textit{Teamfight Tactics}, \textit{Mario Kart},, \textit{Overwatch 2},
\textit{Legends of Runeterra}, \textit{Wild Rift} e outros estão sempre surgindo).
Além disso, contamos com a nossa amada BatIMEduca, bateria do IME juntamente
com a Pedago.

Os dias e horários dos treinos/jogos de cada uma dessas modalidades serão
sempre informados através do instagram e de um mural no vidro da salinha da atlética.
Vocês também podem entrar em contato pelas nossas redes sociais ou diretamente
na página da modalidade (todas estão na nossa bio do Instagram: @atleticaimeusp).

Com essa variedade de modalidades, não tem desculpa pro sedentarismo hein,
bixe? Se você não conhece nenhuma delas, a gente te apresenta e, se já conhece,
vem dar aquele “ooooi, sumido!!!” para aquele esporte que você largou por conta
da Fuvest! Ah, não esqueça de torcer por nosses atletas! Raça e coração, IME!

Além disso, a Atlética promove anualmente campeonatos internos daqueles jogos
que a gente passa hooooras fritando em casa. Já foram promovidos campeonatos de
\textit{Winning Eleven}, \textit{LoL}, \textit{Mario Kart}, \textit{Super Smash Bros}, 
\textit{Mario Tenis}, \textit{Guitar Hero} e 
sinuca.

Ideias e sugestões sobre novas modalidades, campeonatos, inters etc. são
sempre muito bem-vindas! Conversem com a gente!

\end{subsecao}
%%%%%%%%%%%%%%%%%%%%%%%%%%%%%%%%%%%%%%%%%%%%%%%%%%%%%%%%%%%%%%%%%%%%%%%%%%%%%%%%

\begin{subsecao}{Atividades esportivas externas:}

Além das modalidades da AAAMat, você pode participar das atividades do CEPEUSP,
o Centro de Práticas Esportivas da USP! 

O espaço oferece piscina (sim, piscina!!!), campos, quadras e salas incríveis.
Ainda organiza palestras, eventos e dicas sobre nutrição, saúde e bem-estar e 
possui cursos para a Terceira Idade, Comunidade USP e Externa. 

Os cursos vão de alongamento, caminhada e ioga até canoagem, capoeira, futsal, 
basquete, judô, karatê e outros. A duração média é de seis meses, possui custo 
significativo para a matrícula e as inscrições podem ser feitas direto pelo site. 

Para saber mais, visite o site do CEPEUSP: \url{https://www.cepe.usp.br} e acompanhe 
as datas. 

\end{subsecao}
%%%%%%%%%%%%%%%%%%%%%%%%%%%%%%%%%%%%%%%%%%%%%%%%%%%%%%%%%%%%%%%%%%%%%%%%%%%%%%%%

A AAAMat também representa o IME em diversos campeonatos universitários. São
eles:

\begin{subsecao}{BichUSP}

De nome intuitivo e charmoso, o BichUSP é um campeonato disputado entre as
faculdades da USP em que apenas bixes (VOCÊS!) participam. O campeonato
acontece logo nas primeiras semanas de aula, sempre aos finais de semana. Aqui,
vocês têm a chance de suar a camisa IMEana pela primeira vez e ver a torcida
indo ao delírio em cada jogada - ganhando ou perdendo, veteranes estarão
sempre vibrando por vocês!

``Mas, Atlética, eu não sei jogar nenhum desses esportes :('' - Não tem
problema, a gente te ensina! Teremos treinos especiais para que vocês conheçam
a modalidade, es DM’s, es técniques, a gente e outres bixes que te
acompanharão nesse momento único da graduação! O importante aqui é ter vontade
de participar e se divertir!

Fiquem atentos: a Atlética vai divulgar as datas do BichUSP em breve!

Se você acha que tem alergia a esportes, dá uma chance de a gente te mostrar
o contrário! São várias modalidades com várias dinâmicas diferentes, alguma
delas com certeza vai se encaixar no que você gosta! Se quiser só assistir
no começo e vir torcer com a gente, apareça nos jogos e nós vamos gritar
``VERMELHO E BRANCO ATÉ MORRER'' todos juntos!

%REFTIME
%Esse ano o BichUSP acontecerá nas seguintes datas:

%\begin{itemize}
  %\item Tênis todos os dias do BichUSP
  %\item 16 e 17/03/2019 - Basquete e Handebol.
  %\item 23 e 24/03/2019 - Natação, Atletismo, Tênis de Mesa e Xadrez.
  %\item 30/03 e 01/04/2019 - Futebol de campo e Rugby.
  %\item 06 e 07/04/2019 - Vôlei e Futsal.  
%\end{itemize}

Pra vocês se inspirarem, fizemos essa tabelinha que mostra quanto es bixes
brilharam em anos anteriores. Mal podemos esperar pra completar ela com as
conquistas que virão esse ano:

%REFTIME
\begin{center}
  \begin{tabular}{|c|c|}
    \hline
    Ano & Campeão\\
    \hline
    2005 & Basquete Masculino \\
    2005 & Tênis de Mesa Feminino \\
    2007 & Atletismo Masculino\\
    2009 & Atletismo\\
    2011 & Tênis de Campo Feminino\\
    2012 & Basquete Feminino\\
    2014 & Tênis de Mesa Masculino\\
    2014 & Futsal Masculino\\
    2015 & Futsal Masculino\\
    2016 & Vôlei Masculino\\
    2017 & Xadrez\\
    2017 & Rugby Misto\\
    2017 & Rugby Feminino (IME+EEFE)\\
    2018 & Xadrez\\
    2018 & Futebol de Campo Feminino\\
    2022 & Xadrez\\
    2023 & Brawl Stars\\
    2024 & BORA BIXES!!!!\\
    \hline
  \end{tabular}
\end{center}

\end{subsecao}
%%%%%%%%%%%%%%%%%%%%%%%%%%%%%%%%%%%%%%%%%%%%%%%%%%%%%%%%%%%%%%%%%%%%%%%%%%%%%%%%
\begin{subsecao}{Copa USP}

A Copa USP é o primeiro campeonato após o BichUSP, e existem duas séries (Azul:
1ª divisão e Laranja: 2ª divisão). São nesses jogos que colocamos em prática
tudo o que fizemos nos treinos semanais para brilharmos nos jogos da fase de
grupos e então seguirmos arrasando nos jogos mata-matas.

\end{subsecao}
%%%%%%%%%%%%%%%%%%%%%%%%%%%%%%%%%%%%%%%%%%%%%%%%%%%%%%%%%%%%%%%%%%%%%%%%%%%%%%%%
\begin{subsecao}{Jogos da Liga}

Acontece no segundo semestre (UFA, já passou a P1 de cálculo, vem Cálculo 2!!!
\sout{Ou não}) e nessa competição não existe separação por séries. Somente as
faculdades da USP jogam, as disputas são sorteadas para formarem grupos e
então apenas os melhores colocados seguem para a fase final. Essa é a
oportunidade perfeita para gritar um ``CHUPA, POLI!" na arquibancada.

\end{subsecao}
%%%%%%%%%%%%%%%%%%%%%%%%%%%%%%%%%%%%%%%%%%%%%%%%%%%%%%%%%%%%%%%%%%%%%%%%%%%%%%%%
\begin{subsecao}{NDU}

Esse campeonato acontece duas vezes ao ano, e várias faculdades de São Paulo
(tanto da USP quanto algumas não-USP) competem na cidade com rodadas durante
os fins de semana em busca dos melhores
resultados. Confira com ê DM da modalidade se o time está participando da
competição.

\end{subsecao}
%%%%%%%%%%%%%%%%%%%%%%%%%%%%%%%%%%%%%%%%%%%%%%%%%%%%%%%%%%%%%%%%%%%%%%%%%%%%%%%%
\begin{subsecao}{BIFE}

O BIFE É O MELHOR EVENTO ESPORTIVO DO MUNDO! As iniciais das quatro fundadoras
(Bio, IME, FAU e ECA) formam a sigla que dá nome a esses jogos universitários
que a gente tanto ama. Trata-se de um campeonato entre nove faculdades da USP:
a VET, Física, FFLCH, Química, Pedago e, é claro, as quatro fundadoras já 
citadas, às vezes contando também com algumas faculdades convidadas.

Funciona assim: em um determinado feriado, pessoas que jogam, torcem, festejam e
simpatizam se deslocam até alguma cidade do interior do Estado de SP. A cidade nos
dá um alojamento (leia-se: local para tomar um banho quentinho e descansar no
aconchego de sua barraca), alguns ginásios e um local para as festas. São
quatro dias, às vezes três, muito divertidos e engraçados, onde há rivalidade 
apenas dentro de quadra - porque fora é muito amor e integração! Além disso,
o BIFE organiza as três maiores festas da USP: Desmame, Engorda e Abate, as
quais antecedem a viagem para a cidade e servem para todo mundo se divertir e
se conhecer antes do evento principal!

Nosso histórico neste Inter é de parar o trânsito! Olhem só:

%REFTIME
\begin{center}
  \begin{tabular}{|c|c|c|}
   \hline
   Ano & Cidade & Campeão\\
   \hline
   1999 & Jacareí & IME\\
   2000 & Não Houve & - \\
   2001 & Serra Negra & IME\\
   2002 & Socorro & ECA\\
   2003 & São Sebastião & IME\\
   2004 & Cruzeiro & FFLCH\\
   2005 & Jacareí & FFLCH\\
   2006 & Lorena & IME\\
   2007 & Piedade & IME\\
   2008 & Itapeva & IME\\
   2009 & Cruzeiro & IME\\
   2010 & Barra Bonita & IME\\
   2011 & Casa Branca & IME\\
   2012 & Barra Bonita & IME\\
   2013 & Sumaré & ECA\\
   2014 & Cidade/Araraquara & FFLCH\\
   2015 & Taquaritinga & FFLCH\\
   2016 & Registro & FFLCH\\
   2017 & Avaré & FFLCH\\
   2018 & Casa Branca & FAU\\
   2019 & Franca & FFLCH\\
   2020 & On-line & AAAGW\\
   2021 & On-line & IME\\
   2022 & Itapeva & ICBIÓ\\
   2023 & Casa Branca & ICBIÓ\\
   2024 & VAMO IME!! & \\
   \hline
  \end{tabular}
\end{center}

%REFTIME
Em 2020 e 2021, o BIFE se reinventou trazendo o inter de maneira on-line, através do
E-Bife, em que rolaram jogos dos e-sports e também algumas festas. Nas últimas 
duas edições nós ficamos em  4º!!!!! Nossos times contam com vocês para conquistar
mais um título em 2024! VAMO QUE VAMO, GALERA!!!


\end{subsecao}
%%%%%%%%%%%%%%%%%%%%%%%%%%%%%%%%%%%%%%%%%%%%%%%%%%%%%%%%%%%%%%%%%%%%%%%%%%%%%%%%

\begin{subsecao}{Títulos}

Fruto de muito treino, empenho, suor, torcida e amor pelo IME-USP, reunimos
abaixo algumas de nossas conquistas:

%REFTIME
\begin{center}
  \begin{tabular}{|c|c|c|c|}
    \hline
    Ano & Campeonato & Modalidade & Colocação\\
    \hline
    2009 & Intercalouros  & Atletismo       & 1º\\
    2011 & LUPAA          & Atletismo       & 1º\\
    2011 & BOBPAI         & Baseball        & 1º\\
    2016 & BOBPAI         & Baseball        & 2º\\
    2016 & Liga Paulista  & Baseball        & 3º\\
    2017 & Copa USP       & Baseball Fem.   & 2º\\
    2016 & Wakaba         & Softball        & 2º\\
    2016 & Softparty      & Softball        & 1º\\
    2023 & Copa USP       & Softball        & 2º\\
    2024 & Copa USP       & Softball        & 3º\\
    2024 & GECEBS         & Softball        & 1º\\
    2011 & Jogos da Liga  & Basquete Fem.   & 1º\\
    2012 & Copa Camp      & Basquete Fem.   & 2º\\
    2012 & Copa USP       & Basquete Fem.   & 3º\\
    2019 & BIFE           & Basquete Fem.   & 2º\\
    2024 & InterZebras    & Basquete Fem.   & 2º\\
    2024 & Copa USP       & Basquete Fem.   & 3º\\
    2018 & BIFE           & Basquete Masc.  & 1º\\
    2018 & Jogos da Liga  & Basquete Masc.  & 2º\\
    2018 & Jogos da Liga Série Prata & Basquete Masc. & 1º\\
    2019 & BIFE           & Basquete Masc.  & 1º\\
    2022 & BIFE           & Basquete Masc.  & 2º\\
    2023 & NDU            & Basquete Masc.  & 3º\\
    2023 & Copa USP Série Laranja & Basquete Masc. & 1º\\
    2024 & Jogos da Liga  & Basquete Masc.  & 3º\\
    2024 & NDU            & Basquete Masc.  & 3º\\
    2024 & BIFE           & Basquete Masc.  & 1º\\
    2020 & E-BIFE         & Brawl Stars     & 1º\\
    2021 & E-BIFE         & Brawl Stars     & 1º\\
    2020 & E-BIFE         & Clash Royale    & 1º\\
    2021 & E-BIFE         & Clash Royale    & 1º\\
    2021 & E-BIFE         & CSGO            & 3º\\
    2018 & BIFE           & E-Sports        & 1º\\
    2019 & BIFE           & E-Sports        & 2º\\
    2020 & E-BIFE         & E-Sports        & 2º\\
    2021 & E-BIFE         & E-Sports        & 1º\\
    2022 & BIFE           & E-Sports        & 1º\\
    2023 & BIFE           & E-Sports Absoluto  & 1º\\
    2023 & BIFE           & E-Sports Inclusivo & 1º\\
    2015 & Integramix     & Futebol Campo   & 1º\\
    2017 & Copa USP       & Futebol Campo M & 1º\\
    2018 & BIFE           & Futebol Campo F & 2º\\
    2014 & Interfarofa    & Futsal Fem.     & 1º\\
    2015 & NDU            & Futsal Fem.     & 2º\\
    2017 & Copa USP       & Futsal Fem.     & 1º\\
    2017 & Jogos da Liga  & Futsal Fem.     & 2º\\
    2017 & IMEACHECA      & Futsal Fem.     & 1º\\
    2022 & Interatléticas & Futsal Fem.     & 2º\\
    \hline
         \end{tabular}
\end{center}
\begin{center}
\begin{tabular}{|c|c|c|c|}
  \hline
  Ano & Campeonato & Modalidade & Colocação\\
  \hline
    2011 & Copa USP       & Futsal Masc.    & 1º\\
    2012 & Copa Camp      & Futsal Masc.    & 1º\\
    2013 & NDU            & Futsal Masc.    & 1º\\
    2013 & Jogos da Liga  & Futsal Masc.    & 2º\\
    2015 & Integramix     & Futsal Masc.    & 1º\\
    2015 & NDU            & Futsal Masc.    & 2º\\
    2023 & Jogos da Liga Série Prata & Futsal Masc. & 2º\\
    2023 & Interatléticas & Futsal Masc.    & 1º\\
    2023 & BIFE           & Futsal Masc.    & 1º\\
    2024 & NDU            & Futsal Masc.    & 3º\\
    2024 & BIFE           & Futsal Masc.    & 3º\\
    2024 & Copa USP       & Futsal Masc.    & 3º\\
    2015 & Camp. G-4      & Handebol Fem.   & 1º\\
    2015 & Integramix     & Handebol Fem.   & 1º\\
    2016 & CUPA           & Handebol Fem.   & 2º\\
    2016 & Interfarofa    & Handebol Fem    & 1º\\
    2018 & Copa USP       & Handebol Fem.   & 1º\\
    2006 & Jogos da Liga  & Handebol Masc.  & 2º\\
    2009 & Copa USP       & Handebol Masc.  & 1º\\
    2012 & Copa USP       & Handebol Masc.  & 2º\\
    2016 & IMEACHECA      & Handebol Masc.  & 2º\\
    2018 & Jogos da Liga  & Handebol Masc.  & 2º\\
    2022 & BIFE           & Handebol Masc.  & 1º\\
    2024 & BIFE           & Handebol Masc.  & 2º\\
    2024 & CPU Fupe       & Handebol Masc.  & 2º\\
    2023 & Jogos da Liga Série Prata & Handebol Masc. & 2º\\
    2023 & Copa USP Série Laranja & Handebol Masc.  & 1º\\
    2017 & TUES           & Hearthstone     & 2º\\
    2020 & E-BIFE         & Hearthstone     & 1º\\
    2017 & Copa USP       & Jiu-jitsu       & 2º\\
    2024 & BIFE           & Jiu-jitsu Fem.  & 1º\\
    2024 & BIFE           & Jiu-jitsu Masc. & 1º\\
    2022 & BIFE           & Judô Masc.      & 2º\\
    2023 & BIFE           & Judô Masc.      & 2º\\
    2024 & BIFE           & Judô Masc.      & 1º\\
    2021 & E-BIFE         & LOL             & 2º\\
    2018 & BIFE           & Natação Masc.   & 2º\\
    2018 & Jogos da Liga  & Natação Masc.   & 2º\\
    2019 & BIFE           & Natação Masc.   & 3º\\
    2023 & BIFE           & Natação Masc.   & 1º\\
    2024 & BIFE           & Natação Masc.   & 2º\\
    2021 & E-BIFE         & Rocket League   & 2º\\
    2018 & BIFE           & Rugby Masc.     & 2º\\
    2022 & BIFE           & Sinuca          & 1º\\
    2023 & BIFE           & Sinuca          & 1º\\
    2017 & Jogos da Liga  & Tênis Campo M   & 2º\\
    2021 & E-BIFE         & Valorant        & 2º\\
    2016 & IMEACHECA      & Vôlei Fem.      & 1º\\
    2016 & Gran Prix USP  & Vôlei Fem.      & 3º\\
    2017 & Gran Prix USP  & Vôlei Fem.      & 3º\\
    2024 & Quadrangular da EACH & Vôlei Fem. & 3º\\
    2015 & Camp. G-4      & Vôlei Masc.     & 2º\\
    2016 & Copa USP       & Vôlei Masc.     & 1º\\
    \hline
         \end{tabular}
\end{center}
\begin{center}
\begin{tabular}{|c|c|c|c|}
  \hline
  Ano & Campeonato & Modalidade & Colocação\\
  \hline
    2017 & Copa USP       & Vôlei Masc.     & 2º\\
    2017 & NDU            & Vôlei Masc.     & 1º\\
    2018 & Copa USP       & Vôlei Masc.     & 1º\\
    2018 & BIFE           & Vôlei Masc.     & 2º\\
    2018 & NDU            & Vôlei Masc.     & 2º\\
    2024 & Liga USP       & Vôlei Masc.     & 1º\\
    2023 & Jogos da Liga Série Prata & Vôlei Masc. & 2º\\
    2017 & NDU            & Xadrez          & 1º\\
    2017 & Copa USP       & Xadrez          & 3º\\
    2017 & Jogos da Liga  & Xadrez          & 2º\\
    2020 & E-BIFE         & Xadrez          & 1º\\
    2021 & E-BIFE         & Xadrez          & 1º\\
    2022 & BIFE           & Xadrez          & 3º\\
    2023 & NDU            & Xadrez          & 3º\\
    2023 & Copa USP       & Xadrez          & 1º\\
    2023 & BIFE           & Xadrez          & 1º\\
    2024 & BIFE           & Xadrez          & 1º\\
    \hline
  \end{tabular}
\end{center}

\end{subsecao}
%%%%%%%%%%%%%%%%%%%%%%%%%%%%%%%%%%%%%%%%%%%%%%%%%%%%%%%%%%%%%%%%%%%%%%%%%%%%%%%%
\begin{subsecao}{Outras atividades}

\begin{subsubsecao}{Vendas}

A Atlética também quer te ajudar a vestir o vermelho e branco (que, a essa
altura, já corre em suas veias! :D) e traz pra você diversos produtos
personalizados, tais como: canecas, talabartes, corta-ventos,
samba-canção, camisetas e agasalhos do
IME, pra você sair por aí esbanjando seu amor pelo IME-USP $<$3

Vocês podem comprar esses produtos na salinha da Atlética, ou nos nossos
stands de venda.

\end{subsubsecao}
%%%%%%%%%%%%%%%%%%%%%%%%%%%%%%%%%%%%%%%%%%%%%%%%%%%%%%%%%%%%%%%%%%%%%%%%%%%%%%%%
\begin{subsubsecao}{Festas}

%REFTIME
A Atlética e o CAMat já promoveram muitas festas e happy hours. Atualmente,
promovemos a I Will SurvIME, a JunIME, a Melhores do Ano e, além disso, 
auxiliamos as festas pré-BIFE (Desmame, Engorda e Abate). Todas imperdíveis!
Esperamos vocês!

\end{subsubsecao}
%%%%%%%%%%%%%%%%%%%%%%%%%%%%%%%%%%%%%%%%%%%%%%%%%%%%%%%%%%%%%%%%%%%%%%%%%%%%%%%%
\begin{subsubsecao}{Como falar com a Atlética?}

A salinha da AAAMat é a B-18. Ela fica dentro da vivência e é pequena, mas
sempre cabe mais um! Sempre que precisarem conversar com a Atlética, vocês
podem ir até lá e falar com qualquer membro da gestão. Além disso, também temos
outros meios de contato, tais como o e-mail: aaamat.presidente@gmail.com  e as
nossas redes sociais (principalmente o Instagram)

Para ficarem sabendo tudo que acontece na atlética basta:

\begin{itemize}
  \item Acompanhar o site da atlética: \url{https://atletica.ime.usp.br}
  \item Curtir nossa página no Facebook: \url{https://fb.com/aaamat.ime}
  \item Seguir a gente no Instagram: @atleticaimeusp
\end{itemize}

%REFTIME
Talvez vocês notem que existem outros dois instas antigos da AAAMat
(@aaamat\_imeusp e @atleticaime), ambos já não são mais utilizados
e se algum bixe souber como derrubar essas contas, fale para a atlética, 
pois a gestão será eternamente agradecida pela sua ajuda! 

A Atlética inicia 2024 esperando vocês, bixes querides, para que juntos
possamos trazer muitos títulos e troféus para casa!! É importantíssimo que
vocês saibam que estamos abertos para qualquer tipo de crítica, dúvida, ideia
ou sugestão.

Vocês são SEMPRE muito bem-vindes em nossa sala, atividades, times e eventos!
\end{subsubsecao}
\end{subsecao}
\end{secao}
\pagebreak


% Dissecando os Cursos ---------------------------------------------------------
\input{dissecando_os_cursos.tex}
\pagebreak

% AACs e ATPAs ------------------------------------------------------------------
\input{aac_atpa.tex}
\quadrinhos{6}
\pagebreak
% Júpiter Web ------------------------------------------------------------------
\begin{secao}{JupiterWeb}

O \sout{Luciferweb} Jupiterweb é o sistema que administra (quase) tudo que
vocês, bixes, fazem na sua graduação. Desde as suas matrículas até a emissão de
um documento para pagar meia no cinema, o Jupiterweb vai ser muito útil em
vários momentos se vocês conhecerem tudo que ele tem a oferecer. Por isso,
assim que puderem, procurem explorar cada uma de suas funcionalidades.

Talvez você esteja estranhando e pensando "me falaram que é o Portal de Serviços da USP
que faz essas coisas". Tecnicamente, você não está errado: o \url{https://portalservicos.usp.br}
tem exatamente as mesmas funções que o JupiterWeb dentro da aba Graduação.
Entretanto, ele tem mais algumas abas e funções, como a versão da pós do Jupiterweb (o Janus),
tudo condensado em um só site para facilitar o acesso.
Por isso, inclusive, ele também é conhecido como o Olimpo (sério, acessando \url{https://olimpo.usp.br}
ele redireciona para o Portal de Serviços).

Você deve estar se perguntando então: por que nosso guia foca no Jupiter invés do Olimpo?
Basicamente porque os dois portais fazem a mesma coisa e por estamos acostumados com o Júpiter,
que existe há mais tempo que o Olimpo.

\begin{subsecao}{Calendário}

Muito útil pra saber quando é \sout{feriado} que abre a matrícula ou o período
de requerimento.

\end{subsecao}

\begin{subsecao}{Cursos de ingresso}

Aqui, você pode ver as informações sobre o seu curso, o projeto pedagógico e
a grade curricular. Para isso, basta inserir a unidade do seu curso (no caso,
IME) e o nome do seu curso.

Preste atenção na grade curricular, pois ela tem todas as disciplinas que
vocês terão que cursar durante a graduação. Elas estão separadas em obrigatórias
(matérias que terão que cursar para se formar), optativas eletivas (matérias
que terão que escolher dentre algumas opções pra fazer) e optativas livres
(aquelas que acham que você pode se interessar por fazer).

Tome cuidado principalmente com as obrigatórias que são pré-requisitos de
alguma outra matéria! Você não poderá cursar uma matéria se não tiver passado
nos seus respectivos pré-requisitos. Então, \textbf{passem em Cálculo!!!!}

Existem também os requisitos fracos (mas são uma raridade). Neles, vocês não 
precisam ter passado na matéria para cursar a outra, basta ter tirado 3.

\end{subsecao}

\begin{subsecao}{Disciplinas}

Mostra todas as disciplinas da USP (até as que não estão mais disponíveis).
Vocês podem pesquisar por código, por algum trecho do nome ou por unidade. 
Além de ser útil para conferir algumas informações sobre aquela matéria que 
você precisa cursar, ainda ajuda a procurar matérias diferentes pela USP 
(como “Pintura e suas técnicas”, “Psicologia da Morte” ou “Pesca Sustentável”).

\end{subsecao}

\begin{subsecao}{Grade Horária}

Lista as disciplinas em que vocês estão inscritos/matriculados/pendentes, com os
horários de cada uma. É onde geralmente vocês conferem se tá tudo certo com a
sua matrícula ou não. Quando as aulas começam, o ideal é que apareça
“Matriculado” (a matéria fica cinza). Se alguma disciplina aparecer como
pendente (fica azul), então corram na seção de alunos para conferir o que está
errado!

\end{subsecao}
\begin{subsecao}{Histórico Escolar}

O histórico escolar é um documento enorme com tudo de bom (e de ruim) que
aconteceu na sua vida acadêmica na USP. Lá vocês podem conferir cada
disciplina cursada (com nota), as disciplinas que vocês trancaram, seus
aproveitamentos de estudo, quantos créditos vocês já fizeram, a porcentagem de
conclusão do curso e a tão mística “Média ponderada”.

\end{subsecao}

\begin{subsecao}{Acompanhamentos}

\begin{itemize}
  \item \textbf{Evolução no Curso}

    É o jeito visual de acompanhar as disciplinas que vocês já fizeram ou que
    estão faltando. Sua principal função é ser printado ao final do curso (com
    tudo verde) para postar no Facebook e fazer sucesso com a família na
    internet.
  \item \textbf{Rendimento Acadêmico}

    Mais um lugar pra vocês conferirem suas notas, quantidade de matérias
    cursadas, e quantidade de créditos cursados, só que aqui tudo aparece em
    forma de gráfico.
  \item \textbf{Perfil de deficiências}

    Questionário da USP sobre deficiências.

  \item \textbf{Meus benefícios e bolsas}

    Onde ficam registradas as suas bolsas e alguns serviços prestados na
    universidade que valem dinheirinhos. Também mostra as datas que a USP lhes
    paga para vocês organizarem sua vida financeira!!

  \item \textbf{Dados pessoais}

    O nome é bem intuitivo. Mantenham tudo atualizado para não ter problemas com
    a SPTrans, SAS ou qualquer coisa que confira seus dados em algum momento.

  \item \textbf{Bilhete Único SPTrans}
    Onde se solicita o envio dos dados do item anterior para a SPTrans, só assim
    você pode solicitar (no site da SPTrans) o seu bilhete único de estudante.

\end{itemize}

\end{subsecao}

\begin{subsecao}{Cartão USP}

Aqui é onde se solicita o Cartão USP e o BUSP. Cartão USP é o que chamamos de
carteirinha, ela serve para entrar no bandejão, no CEPE, pagar meia, entre outras
coisas, é imprescindível que você a solicite. Já o BUSP é o ``Bilhete Único'' da USP,
com ele você pode utilizar os 3 circulares de graça, então solicite esse também!

\end{subsecao}

\begin{subsecao}{Emissão de Documentos}

Emite alguns documentos que precisamos no dia a dia, como um comprovante de
matrícula ativa. Dependendo do documento, vocês vão precisar ir até a seção de
alunos, mas os mais requisitados vocês encontram por aqui.

\end{subsecao}

\begin{subsecao}{Matrícula}

Aqui começa uma das partes mais legais da Universidade: A matrícula! Pois é,
bixes, vocês não vão cursar apenas o que a sua instituição de ensino manda, mas o
que vocês quiserem! Apesar de ter uma grade de matérias obrigatórias, vocês
podem cursá-las na ordem que lhes for mais conveniente e sempre ir intercalando
com optativas eletivas e optativas livres.

Não que isso seja uma boa prática. O seu curso foi pensado e montado com uma
sequência de matérias que faça sentido na sua vida, mas vocês têm independência
para quebrar isso caso seja melhor pra vocês.

\textbf{Vocês não precisam se preocupar com isso agora! As matérias do primeiro
semestre já são inseridas automaticamente, vocês só precisam fazer a própria
matrícula a partir do segundo semestre.}

Quando a hora chegar, vocês vão clicar em “Fazer matrícula” e escolher o tipo de
matéria que vocês querem cursar: Obrigatórias, Optativas eletivas, Optativas
Livres ou Extracurriculares. Quando vocês escolhem um desses blocos, vão
aparecer uma série de matérias disponíveis pra vocês com as turmas (cada turma
tem um horário e um professor específicos). Basta vocês irem escolhendo quais
querem cursar, conferir umas 15 vezes e salvar. \textbf{Não se esqueçam de salvar,
bixes, ou não estarão matriculados nas disciplinas!}

Pode até parecer um monstro de 7 cabeças, mas está beeeeem longe disso!
Aproveitem esse momento para consultar aquele veterane que você ama muito e que
vai te mostrar como é tudo bem tranquilo.

Mas quando é mesmo que eu faço isso?

O Jupiterweb vai abrir algumas “Interações de matrícula”, geralmente duas ou três.
Apenas nesses dias, o botãozinho de “Fazer matrícula” fica disponível. Vocês
recebem e-mails da USP, fica um banner no meio do saguão do bloco B, fazem
evento no Facebook e divulgam a data no calendário da USP… Não tem desculpa pra
não saber quando é!

No dia de abertura das interações, o Jupiterweb revela suas verdadeiras garras
(e lentidões) dando absolutamente todos os erros possíveis e tentando
atrapalhar a sua vida. Então, talvez esse não seja o melhor horário pra tentar
se matricular. A ordem de matrícula não importa, então podem ficar tranquilos e
fazer em qualquer momento de qualquer uma das interações (MAS NÃO SE ESQUEÇAM DE
FAZER EM ALGUM MOMENTO!!!).

Depois de cada interação, o Jupiterweb vai colocar vocês em um ranking que olha seu
semestre, o que vocês deveriam estar cursando, suas notas e o que mais for
conveniente para decidir se vocês ganham ou não uma vaga na disciplina.
Geralmente isso não é um problema com matérias obrigatórias (principalmente se
vocês seguirem sua grade ideal direitinho) ou matérias com poucos inscritos. Mas
se vocês estiverem tentando uma optativa muito concorrida (como ``Psicologia da
morte''), se preparem que vocês podem perder uma vaga muito sonhada.
Esse resultado pode ser acessado no ``Resultado das consolidações''.

\end{subsecao}
\begin{subsecao}{Requerimentos}

Por vários motivos, vocês podem ter sua matrícula de uma disciplina negada ou
até mesmo nem conseguir se matricular porque essa disciplina não tem vagas
reservadas para o seu curso. Isso não quer dizer que vocês não vão conseguir
cursar, mas que vocês vão ter que enfrentar mais burocracias para isso. Os
requerimentos, também conhecidos como melhores amigos do aluno da graduação,
são pedidos formais em que vocês inserem a matéria da qual vocês foram
chutados/negados/impedidos de cursar e uma boa justificativa. Esse requerimento
vai para o coordenador do seu curso e possivelmente para o professor de
matéria, pessoa que vai decidir se o seu motivo é bom o suficiente e se o
número de vagas comporta a sua presença.

Vocês podem fazer um requerimento para cursar qualquer matéria, por qualquer
motivo, mas se o pedido for muito esdrúxulo, vocês podem receber um não.

O resultado dos requerimentos demora 84 anos e um pouco para sair. Por isso,
vocês vão acompanhar as aulas da disciplina requerida mesmo sem saber o
resultado, confiantes de que vocês vão receber um OK depois de um tempo. Existe
um prazo limite para aprovar ou negar um requerimento, então fiquem tranquilos
que até a P1 ou P2, vocês já devem saber se estão matriculados ou não.

Dica: Conversar com os coordenadores e o professor responsável da disciplina
pode ajudar! Se vocês querem muito cursar uma matéria, não tenham vergonha de
lutar por ela :)

\end{subsecao}

\begin{subsecao}{Vacinação Covid-19} %REFTIME %PANDEMIA

Aqui é o lugar para vocês colocarem os comprovantes de vacina contra a Covid-19! 
O esquema de adionar seus comprovantes é bem tranquilo, mas se rolar alguma
dúvida, já está disponível na mesma página um tutorial de como prosseguir. E
aqui fica o aviso, bixes, NÃO se esqueçam de anexar os comprovantes!! Sem eles
vocês não poderão acessar os prédios da USP, e essa regra vai desde o IME até o
CEPE! Não se preocupem se não tomaram a terceira dose ainda, podem anexar as duas
primeiras e adicionar depois o comprovante da terceira! O importante é não esquecer!

Não fiquem ansiosos se demorar muito para validarem seus dados. Cada um deles está
sendo conferido individualmente. Mas quer uma dica? Baixa o app do E-Card e dá uma
olhada se aparece alguma atualização logo abaixo do seu rostinho, podendo ser
"imunização completa" para quem já possui as três doses ou "imunizado com duas doses".

\end{subsecao}

\end{secao}

\pagebreak

% Um Pouco Sobre o IME ---------------------------------------------------------
\begin{secao}{Um Pouco Sobre o IME}
	
O IME, Instituto de Matemática e Estatística (não, bixes, não está errado.
Computação não faz parte da sigla mesmo! Aliás, faz sim, mas só os inteligentes
podem ver!), tem vários blocos: bloco A, bloco B, bloco C, bloco D e o
novo bloco, até então um mistério para todos. Leia atentamente, porque com essa
confusão toda, até alguns veteranes insistem em dizer que só existem os blocos
A, B e C.


\begin{subsecao}{Bloco A}

No bloco A, as coisas mais importantes são: as salas dos professores, a parte
administrativa (diretoria, secretarias de departamento), algumas salas de aula
(geralmente usadas para a pós-graduação), a IME Jr, a biblioteca, as ET's e a Rede Linux.

\begin{itemize}

\item {\bf IME Jr:} é a empresa júnior do IME, que é administrada pelos próprios
alunos da graduação. Fica na sala 258A.

\item {\bf Biblioteca:} fica logo na entrada, do lado direito. Há algumas mesas
individuais e sofás (ambos muito confortáveis para tirar um cochilo) e algumas salas
com lousa para estudo em grupo. Tem também uma sala de estudos enorme, 
mas lá não pode fazer barulho (cuidado até mesmo ao abrir e fechar a porta!).
Veja mais informações sobre funcionamento e empréstimos adiante.

\item {\bf ET's (Estações de Trabalho):} são salas com alguns terminais que têm
acesso à Internet e muito mais. Não, bixes, não se animem, pois ela não é para
vocês! Essa sala só pode ser usada pelos alunos da pós-graduação e por alguns
alunos que fazem iniciação científica.

\item {\bf Rede Linux:} rede de computadores que podem ser usados por alunos
da graduação (ieiii!). Utilizam o Linux, que, ao contrário do Windows, é um
Sistema Operacional. Todo aluno que fizer uma conta, possui 75 cotas de impressão por mês 
(imprimir de graça).

\end{itemize}

\end{subsecao}

\begin{subsecao}{Bloco B}


No bloco B, vocês devem conhecer as salas de aula, a Vivência, o CEC, 
a seção de alunos, a lanchonete, o CAEM e a Gráfica para vocês 
copiarem \sout{o caderno do colega} o capítulo do livro para estudar para a prova.

\begin{itemize}
\item {\bf Vivência:} sala 18, arrasadora de graduações, é onde as pessoas podem
dormir, jogar sinuca, cartas, xadrez, vídeo-game e até
mesmo estudar. Guardem bem esse nome, porque vocês vão passar a maior parte do tempo lá.
Lá também estão os armários e as salas do CAMat e da Atlética.

Os armários são renovados semestralmente (mas vocês também podem alugar por um ano!). 
Após o período de renovação, os que estiverem sobrando ficam livres para serem alugados. 
Aí cabe a vocês, bixes, ficarem atentos às datas divulgadas pelo CAMat e demonstrarem 
interesse pelos armários para conseguirem entrar no sorteio.

\item {\bf Lanchonete:} uma alternativa para os dias chuvosos em que você está sem
  guarda-chuva para ir até o bandejão é a lanchonete do IME. Tem um vasto
  cardápio de salgados e congelados esquentados no microondas, além de sucos,
  refrigerantes, sorvetes, chocolates, bolachas e similares.

\item {\bf CEC (Centro de Ensino de Computação):} é um centro munido de dezenas
de computadores rodando Windows, e alguns rodando Linux. Também é possível fazer
(ou pelo menos tentar fazer) EPs em situações de desespero (vocês vão descobrir o que é
isso logo, logo). Algumas aulas de computação são ministradas lá (como Desafios
de Programação do BCC).

\item {\bf Seção de Alunos:} é aqui que vocês vão resolver quase todos os
pepinos de vocês durante a Graduação. Desde descobrir por que o Júpiterweb não
quer matricular vocês em uma disciplina obrigatória até \sout{trancar essa mesma
disciplina mais ou menos um mês depois} pegar o BUSP provisório e o definitivo. %REFTIME (se em algum momento a matrícula presencial voltar a ser na seção de alunos, substituir o trecho sobre trancamento/BUSP pela data da matrícula)

A Seção de Alunos fica na sala 12, ao lado do CEC.
\begin{itemize}
\item[-] Horário de atendimento presencial: segunda, quarta e sexta-feira, das 11h30 às 13h e das 19h às 20h30.
\item[-] E-mail: \tt{graduacao@ime.usp.br}
\item[-] Telefone: \tt{11 3091-6149} e \tt{11 3091-6279}
\end{itemize}

Neste ano, a Seção de Alunos nos informou
que o quadro de funcionários está bem reduzido, por isso, pedimos que:
\begin{itemize}
\item[-] Deem preferência ao uso do Sistema JupiterWeb (troquem conversas, ideias e informações com os colegas - um pode ajudar o outro)
\item[-] Evitem encaminhar mensagens por e-mail, pois devido a grande demanda de serviço, o retorno tem sido demorado.
\item[-] Evitem as ligações por telefone, pois a prioridade tem sido o andamento do serviço interno. 
\end{itemize}

\item {\bf CAEM:} sigla para o Centro de Aperfeiçoamento do Ensino da
  Matemática. É um órgão de extensão, que oferece cursos, oficinas, palestras e
  presta serviços de assessoria para professores de Matemática. Os que fazem
  Licenciatura (futuros professores) também podem usufruir dos serviços do CAEM,
  e até mesmo estagiarem lá. Fica no 1º andar do bloco B, em frente às escadas.

\item{\bf Gráfica:} fica no térreo, saindo da vivência e indo até o final do segundo
  corredor à direita (o que não tem os banheiros no final). Lá vocês podem fazer cópias
  e imprimir colorido, mas, para isso, precisam comprar fichas na salinha da Atlética 
  (que fica na Vivência) ou então na Tesouraria(no bloco A).
\end{itemize}

\end{subsecao}

\mapaime{bloco_A_terreo}

\mapaime{bloco_A_piso1}

\mapaime{bloco_A_piso2}

\mapaime{bloco_B_terreo}

\mapaime{bloco_B_superior}

\begin{subsecao}{Bloco C}

Apesar de atualmente abrigar os professores e a secretaria do departamento de
Computação, o misterioso Bloco C é um local a que nós, estudantes, infelizmente não
temos livre acesso. Para garantir tranquilidade e boas condições de trabalho aos
docentes, os alunos devem ser anunciados ou apresentarem uma boa desculpa para
adentrar o local. (Na verdade, os professores do MAC acham que os alunos são monstros
verdes e gosmentos, com $e^{10}$ braços, $\pi$ olhos e que comem criancinhas;
por isso, não querem esse tipo de criatura perambulando pelos seus corredores.)

\end{subsecao}

\begin{subsecao}{Bloco D (vulgo C')}

O bloco D não passa de uma extensão do bloco C. Na verdade, eles são o mesmo bloco,
só que têm entradas independentes. Pertence ao NUMEC, que ninguém sabe ao certo
o que significa. Alguns dizem que ele não existe, e que é apenas fruto da sua
imaginação.

\end{subsecao}

\begin{subsecao}{CCSL}

O anexo do Bloco C abriga o CCSL --- Centro de Competência em \textit{Software Livre}.
O que tem lá? Algumas salas de docentes (afinal, ele é uma extensão do
Bloco C), alguns laboratórios (de Sistemas, de Inteligência Artificial,
de Visão Computacional, de Computação Musical...), um pequeno auditório
onde rolam palestras e seminários e, finalmente, o laboratório de extensão,
que é usado para cursos rápidos e palestras com parte prática em que público
precise de um computador. É esse laboratório, que abriga grande parte dos grupos
de extensão do IME, que vocês verão na seção ``Atitude, bixes!''.

Se vocês se interessam por \textit{Software Livre}, não deixem de ler a seção
do CCSL desse Guia!

\end{subsecao}

\begin{subsecao}{Biblioteca do IME-USP}

A Biblioteca Carlos Benjamin de Lyra é uma das 48 bibliotecas da Agência de Bibliotecas e Coleções Digitais da USP (ABCD). Considerada uma das maiores bibliotecas na área de matemática da América Latina, seu acervo é formado por livros, relatórios, revistas acadêmicas, teses, dissertações, trabalhos de docentes, CDs e DVDs, totalizando cerca de 235.000 itens. 

\begin{subsubsecao}{Instalações:}

O espaço físico da biblioteca é de 1.642m\textsuperscript{2}. Possui um amplo salão de estudos com mesas individuais, 13 salas para estudo em grupo e bancadas para estudo individual, totalizando 221 assentos.
\end{subsubsecao}

\begin{subsubsecao}{Empréstimos:}

O aluno de graduação poderá pedir emprestado 10 livros por 10 dias. Para o empréstimo é necessário cadastrar uma senha pessoalmente e apresentar qualquer um dos seguintes documentos: cartão provisório, carteirinha USP, aplicativo e-card ou RG/CNH. Poderão ser feitas 3 renovações pela internet (DEDALUS) e 3 reservas de livros que estiverem emprestados com outros alunos. O atraso na devolução gera suspensão (nº livros atrasados multiplicado pelo nº dias em atraso). 
\end{subsubsecao}

\pagebreak

\begin{subsubsecao}{Recursos para pesquisas:}
\vspace{-15pt}
\begin{itemize}
    \item DEDALUS - é o catálogo online para consulta e localização dos materiais nos acervos das bibliotecas da USP e para renovação de empréstimos e pedidos de reservas. 
    \item Portal de Busca Integrada - interface única para pesquisa de materiais impressos das bibliotecas e dos documentos disponíveis online, possibilitando acesso a textos completos. 
    \item Bases de Dados e Portal Capes – acesso online a e-books e artigos de revistas assinadas pela USP e pela Capes.
    \end{itemize}
\end{subsubsecao}

\begin{subsubsecao}{Serviços oferecidos:}

Além dos empréstimos, a Biblioteca possui scanner de autoatendimento e oferece os seguintes serviços para os alunos:
\vspace{-15pt}
\begin{itemize}
    \item COMUT (comutação bibliográfica): obtenção de cópias de artigos e capítulos de livros nas principais bibliotecas brasileiras; 
    \item EEB (empréstimo entre bibliotecas): empréstimo entre bibliotecas de instituições externas à USP ou bibliotecas USP interior;
    \item Orientação para trabalhos acadêmicos: orientação para formatação de TCC, teses, dissertações, citações e referências bibliográficas;
    \item Treinamentos e workshops: uso dos recursos da biblioteca e pesquisas em portais e bases de dados;
    \item Apoio ao pesquisador: auxílio na localização de artigos, livros e referências bibliográficas;
    \item Aplicativo Bibliotecas USP: pesquisa no acervo das bibliotecas da USP diretamente do Android, iPhone, iPod Touch ou iPad; 
    \item Conexão VPN: acesso remoto (de casa) às bases de dados, revistas eletrônicas e e-books. 
\end{itemize}
\end{subsubsecao}

\begin{subsubsecao}{Horário de funcionamento:}

Período letivo: de 2ª a 6ª feira das 8h00 às 21h30 \newline
%\phantom{Período letivo: }de sábados das 9h00 às 13h00 \newline%
Período não letivo: de 2ª a 6ª feira das 8h00 às 20h00 \newline
%\phantom{Período não letivo: }de sábados fechada%
O salão de estudos tem uma entrada independente e fica aberto 24h para aqueles que querem 
virar a noite estudando, mas apenas estudantes da comunidade IME podem permanecer lá entre as 23h30 e as 6h.
\end{subsubsecao}
\begin{description}
  \item[Instagram:] \url{https://www.instagram.com/bibliotecaimeusp/}
  \item[Facebook:] \url{https://www.facebook.com/bibliotecaimeusp/}
  \item[Site:] \url{https://www.ime.usp.br/bib}
  \item[E-mail:] \texttt{bib@ime.usp.br}
  \item[Telefone:] 3091-6109 e 3091-6174
  \item[Whatsapp:] 3091-6174
\end{description}

\end{subsecao}
\end{secao}


% Atitude, bixo! ---------------------------------------------------------------

\begin{secao}{Atitude, bixes!}
\quadrinhos{7}

Na USP, alunes têm a liberdade e apoio de se organizarem
para montar grupos de debates, ciclos de palestras, grupos
de desenvolvimento e até mesmo grupos para jogarem alguma coisa (como
RPG, Magic, Yu-Gi-Oh ou algum esporte).
Portanto, caso vocês tenham algum projeto em mente, não hesitem
em se organizarem com seus amigos e se informarem sobre como avançar com essa
ideia. Lembre-se de que seus veteranes estão aí para te aconselhar e tirar
eventuais dúvidas.

É possível também juntar-se com alguns amigos e formar grupos de
estudo, seja para alguma matéria com a qual vocês tenham dificuldade, para
discutir aquele EP/lista de exercícios que ninguém está conseguindo
fazer ou simplesmente estudar algum tópico de interesse mútuo.

A seguir, temos (majoritariamente) alguns dos exemplos dos grupos que foram 
direta ou indiretamente criados por alunes do nosso Instituto!

% RDs --------------------------------------------------------------------------
\begin{subsecao}{RDs}

Antes de mais nada, RD significa Representante Discente. O RD é o estudante 
que representa nossos interesses frente aos diversos conselhos e comissões 
institucionais existentes, sendo por vezes um elo de ligação entre professores 
e alunos. O RD atua na tomada de decisões internas do IME: mudanças no currículo, 
aumento de vagas na FUVEST e distribuição das cotas, reformas, mudança no corpo 
docente, enfim, coisas desse tipo.

Dá para perceber o quanto é importante ter um estudante em cada um desses conselhos 
visando o repasse e a transparência aos estudantes de questões administrativas do 
instituto. Então, ao assumir esse cargo, vocês poderão entender melhor o 
funcionamento do IME e, organizar coletivamente com o centro acadêmico, o processo de 
melhorá-lo.

Desde 2017, as eleições são organizadas anualmente pela administração do instituto 
ao final do segundo semestre. Anteriormente, as eleições eram organizadas pelo centro 
acadêmico, que criava uma Comissão Eleitoral independente e composta por membros do 
centro acadêmico e demais estudantes do IME interessados em organizar o processo eleitoral. 
Essa alteração permitiu a adoção de eleições on-line por meio do Helios Voting, e, embora 
haja o argumento de que os métodos virtuais sejam mais participativos, o quórum de votação 
nas eleições para RDs é quase metade do quórum das eleições das entidades representativas.

Bom, agora vejamos um breve resumo do que mais ou menos acontece em cada um dos 
colegiados nos quais temos direito a representante(s):

Existem diferentes níveis de hierarquia na administração.

{\bf As CoCs,
Comissões Coordenadoras de Curso (Lic, Pura, Estatística, Aplicada, BMAC e
Computação)} são as mais próximas dos alunos. Temos um cargo de aluno em cada
comissão. São comissões pequenas, que tratam dos problemas internos de cada
curso: mudança de currículo, requerimentos, optativas etc. São subordinadas
à CG e ao conselho do relativo departamento. Analogamente, temos um cargo em cada
Comissão Coordenadora de Programa (de Pós).

{\bf Os Conselhos de Departamento (MAT, MAE, MAC e MAP)} têm uma dinâmica um
pouco diferente das CoCs: são mais formais. Cada conselho se reúne (quase)
mensalmente e são formados (em geral) por mais pessoas, sendo que existem
regras sobre participação dos diferentes níveis hierárquicos de
professores (Titular, Associado, Doutor e Assistente). Nesses conselhos, além
de aprovar algumas das decisões das Comissões Coordenadoras de Curso e de
Programa (pós) e distribuição de carga didática, são discutidos reoferecimento
de curso, revisão de prova, supervisão das atividades dos docentes,
afastamentos (temporários ou não), contratação de professores e muitas outras
coisas.
Os Conselhos de Departamento são subordinados à Congregação e ao CTA.

{\bf A Comissão de Graduação (CG)}, basicamente, avalia requerimentos,
mudança/criação de cursos e jubilamentos. Analogamente, existe a Comissão de
Pós-Graduação (CPG). Ambas são subordinadas à Congregação.

{\bf A Comissão de Inclusão e Pertencimento (CIP)} lida com a questão de Permanência,
acompanhando a distribuição de auxílios feita pela reitoria, além de acompanhar
questões de saúde mental e eventualmente receber denúncias sobre opressões e 
violências que possam acontecer no instituto.

{\bf A Comissão de Cultura e Extensão (CCEx)} Cuida das atividades de extensão: 
Matemateca, CAEM, AACs que sejam extensão, etc.

Também há comissões mais específicas, como a comissão de estágio, a comissão de
pesquisa (do doutorado) e o Centro de Competência em Software Livre (CCSL), da
computação.

Os dois conselhos mais importantes são o CTA e a Congregação, ambos presididos
pelo Diretor.

{\bf O Conselho Técnico e Administrativo (CTA)} cuida de todas as questões não
acadêmicas: Orçamento, reformas, avaliação dos funcionários, Xerox etc. É
formado pelos quatro chefes de departamento, diretor, vice-diretor, um
representante dos funcionários e um RD.

{\bf A Congregação} é o órgão máximo do Instituto. Inclui muitos professores, a
maioria titular. São 3 RDs de graduação e 2 de Pós. Basicamente,
nesse órgão, são rediscutidas e aprovadas (ou não) muitas das decisões
dos órgãos subordinados. Os membros da Congregação têm voto na eleição para
Reitor e Vice-Reitor.

Bom, caso não tenha ficado claro desde o começo desse texto, percebam que é
muito importante ter um aluno em cada um desses conselhos. Se estiverem tendo
problemas com professores, requerimentos etc., ou simplesmente quiserem saber
o que anda acontecendo, procurem o RD certo pra conversar. Perguntem,
participem, votem e façam o IME um lugar melhor.

Sobre a eleição dos RDs: A eleição oficial para os RDs acontece no final do ano
(então fiquem atentos!) e é organizada pela diretoria do instituto. Os
interessados devem preencher um formulário de inscrição e levar até a
Assistência Acadêmica do IME, que vai organizar todos os inscritos e abrir um
processo de eleição on-line (em que todo IMEano pode votar). É indicado conversar
com os RDs anteriores para evitar que duas pessoas se inscrevam na mesma vaga e
alguma outra fique sem representação(se organizar direitinho, ocupamos todas as vagas
disponíveis). Quando a eleição
estiver se aproximando, vocês receberão (vários) e-mails com os documentos
necessários, prazos e links de votação.

%REFTIME
O resultado da eleição anterior com os RDs de 2025 pode ser encontrada no site
do CAMat:
\url{https://camat.ime.usp.br/}


\end{subsecao}


% Rede Linux -------------------------------------------------------------------
\input{rede_linux.tex}

% FLUSP ------------------------------------------------------------------------
\input{flusp.tex}

% IME Júnior -------------------------------------------------------------------
\input{ime_jr.tex}

% IME Finance ------------------------------------------------------------------
\input{ime_finance.tex}

% MaratonUSP -------------------------------------------------------------------
\begin{subsecao}{MaratonUSP}

\figurapequenainlineapertada{maratonusp}

O MaratonUSP é o grupo de programação competitiva da USP que tem como foco o
preparo de estudantes para a Maratona de Programação e Olimpíadas de Informática.

No início do ano, o grupo promove o BixeCamp que são aulas direcionadas aos bixes
para o estudo de algoritmos e estruturas de dados, por isso, não deixe de
participar delas. Assim, o MaratonUSP é, atualmente, referência para o
Brasil inteiro por suas diversas conquistas em competições e pelo seu canal do
YouTube, com aulas assistidas por milhares estudantes de programação competitiva
pelo país.

O ambiente de cooperação faz com que o grupo tenha resultados incríveis na
Maratona, sendo Tetracampeão da fase brasileira e tendo participado 16
vezes da etapa mundial da competição, em países como a China, Tailândia,
Rússia, Estados Unidos e muito mais - tudo de graça! Vale lembrar que a própria
fase nacional acontece em um local diferente do país todo ano, o que também gera
uma ótima oportunidade para conhecer mais lugares do Brasil.

Por fim, a vida de Maratoneire não impacta só sua passagem pela universidade,
ela também abre portas para o futuro. É bastante comum que ex-participantes
trabalhem em empresas como Google, Meta e Microsoft, tanto no Brasil,
quanto no exterior. As habilidades desenvolvidas durante as competições,
viajando pelo mundo, dando aulas, preparando vídeos para o YouTube ou
simplesmente estudando com colegas durante uma tarde serão úteis tanto no
mercado de trabalho quanto na academia, além de render anos de muita diversão!

\begin{description}
\item[Facebook:] \url{https://facebook.com/MaratonUSP}
\item[Site:] \url{https://www.ime.usp.br/~maratona}
\item[Youtube:] \url{https://www.youtube.com/c/maratonusp}
% Atualizar o link do Telegram todo ano (esse é específico de 2024)
\item[Telegram:] \url{https://t.me/bixes2025}
\end{description}

\end{subsecao}


% USPCodeLab -------------------------------------------------------------------
\begin{subsecao}{USPCodeLab}

\figurapequenainlineapertada{uspcodelab}

O USPCodeLab é um grupo de extensão que tem por objetivo criar um espaço
colaborativo para criar e incentivar o desenvolvimento de tecnologia na USP.

Nosso foco é aprender na prática ferramentas e técnicas de desenvolvimento de software que permitam
solucionar problemas do mundo real.

Durante o semestre, organizamos o webdev que é um curso dado por membros do codelab para ensinar o
básico de web (Html, Css e Javascript). Ao ganhar mais confiança no desenvolvimento web formamos
grupos de estudos para desenvolver projetos legais propostos pelos próprios participantes (básicos
e avançados). Alguns exemplos são um sistema de reserva de armários do CAMat e o grupo devboost que
desenvolveram um site para cadastrar oportunidades (empregos, IC, \dots). 

O USPCodeLab também organiza hackatons como o shehacks e o hackfools, estes são competições de
programação que os participantes se dividem em grupos para pensar na solução de um problema e tentar
elaborar.

Curtam nossa página do Facebook e entrem no nosso grupo do Telegram para saber
datas e horários das nossas reuniões abertas. Participem do USPCodeLab!

\begin{description}
\item[Facebook:] \url{https://uclab.xyz/facebook}
\item[Telegram:] \url{https://t.me/codelabbutanta}
\item[Site:] \url{https://uclab.xyz/site}
\item[Instagram:] @uspcodelab
\end{description}

\end{subsecao}


% USPGameDev: Pesquisa e Desenvolvimentos de Jogos na USP ----------------------
\input{usp_game_dev.tex}

% IMEsec -----------------------------------------------------------------------
\filbreak
\begin{subsecao}{IMEsec}

\figurapequenainline{imesec}

O IMEsec é um grupo de extensão focado em aprender, estudar e se divertir com a
segurança da informação. Sem que a maioria das pessoas se dê conta, este nicho
está presente no cotidiano. Mandar uma mensagem no WhatsApp, navegar pela web,
entrar no Facebook: exemplos de ações simples do dia-a-dia que necessitam ser
feitas de maneira segura, visando à privacidade do usuário.

Nosso grupo, formado majoritariamente por alunos do IME-USP no início de 2017,
busca entender melhor este universo e expandi-lo no ambiente universitário. O
foco tem sido amplo; desde resolução de desafios on-line (que são muito
divertidos), participação em competições (sim, competições -- \textit{capture the
flag} — muito, muito legais), apresentação de palestras e até desenvolvimento de
projetos que possam beneficiar a população. Todos são bem-vindos; basta ter
interesse pelo assunto.

As experiências decorrentes de nossas atividades ajudaram e ajudam nos estudos,
além de agregarem valor à graduação. Por meio desses, mergulhamos não só em
computação, mas também em matemática, estatística e até outras áreas bem
inusitadas.

Fiquem de olho no nosso grupo do Telegram, ou no nosso servidor do Discord, 
para saber mais sobre as reuniões semanais, além de fatos interessantes (ou 
simplesmente engraçados) sobre o vasto mundo da segurança. Participem!

\begin{description}
  \item[Discord:] \url{https://discord.gg/ZcMvXStFVS}
  \item[Telegram:] \url{https://tiny.cc/imesec-telegram}
  \item[Site:] \url{https://imesec.ime.usp.br/}
\end{description}

\end{subsecao}


% L.E.A.R.N. -----------------------------------------------------------------------
\input{learn.tex}

% SymComp -----------------------------------------------------------------------
\input{symcomp.tex}

% Hardware Livre ---------------------------------------------------------------
\input{hardware_livre.tex}

% Tecs: Computação Social ------------------------------------------------------

\begin{subsecao}{Tecs}

\figurapequenainline{tecs}

Nós somos um grupo de extensão focado no impacto social da computação e da 
tecnologia e que desenvolve projetos em três frentes: educação, ética e 
serviços. Nossos projetos visam, respectivamente, promover a educação tecnológica 
igualitária da população por meio de cursos, oficinas e ações promovidas pelo 
grupo; unir esforços para formar uma sociedade e profissionais éticos e
conscientes sobre o uso da tecnologia; e estimular alunos a usarem a tecnologia 
para solucionar problemas da comunidade local.

Visamos a um cenário de equidade dos saberes, no qual as universidades superem
as barreiras de restrição de conhecimentos e técnicas. Pretendemos, a longo prazo,
contribuir para uma formação universitária que estimule maior consciência social, 
capaz de gerar profissionais da área de computação hábeis em refletir sobre as 
implicações éticas e sociais do seu trabalho, desmentindo, assim, o mito da 
neutralidade tecnológica.

Em termos gerais, pretendemos que os estudantes entendam como a tecnologia pode
ser utilizada para o bem coletivo, e utilizem esse conhecimento na prática, por
meio de colaborações com a comunidade local, os serviços públicos, as
organizações não-governamentais e as sem fins econômicos. 

Se você tem interesse em promover o ensino de computação, em estudar e debater 
questões éticas e sociais no contexto tecnológico, ou em desenvolver aplicativos,
sites ou sistemas em parcerias com projetos sociais, entre em contato conosco e 
participe do grupo!

\vspace{-1em}
\begin{description}
  \item[Site:] \url{www.ime.usp.br/~tecs}
  \item[Facebook:] \url{www.facebook.com/tecs.usp}
  \item[Telegram:] \url{https://t.me/tecsusp}
  \item[Twitter:] \url{https://twitter.com/tecsusp}
  \item[Instagram:] \url{https://www.instagram.com/tecs.usp} 
\end{description}

\end{subsecao}


% DynaUSP: Empreendedorismo
\input{dynausp.tex}

% Diversime --------------------------------------------------------------------

\input{diversime.tex}

% Existimos! -------------------------------------------------------------------
\input{existimos.tex}

% Comissão de Acolhimento da Mulher! -------------------------------------------
\input{cam.tex}

% Comissão de Inclusão e Pertencimento -----------------------------------------
\begin{subsecao}{Comissão de Inclusão e Pertencimento- CAM}

O que é a CIP?

A Comissão de Inclusão e Pertencimento do IME-USP é composta 
por pessoas de diferentes setores do IME e tem por objetivo promover e 
zelar por um ambiente que valorize e acolha a pluralidade através de atividades 
relacionadas a inclusão, pertencimento, diversidade e equidade no âmbito do IME.

Quando procurar a CIP ❓

- Solicitar adaptações pedagógicas ou funcionais para pessoas com TEA e/ou outras condições neurodivergentes.
- Buscar ajuda ligada a questões de saúde mental, seja para você ou para outra pessoa do IME.
- Buscar acolhimento ou denunciar comentários preconceituosos e atitudes inadequadas nos espaços do IME.
- Procure a CIP mesmo em caso de dúvida. Se a sua demanda não for de nossa competência, podemos ajudar a encaminhar para outras instâncias.

Onde nos encontrar?

Atendimento presencial:
Sala 260, Bloco A do IME-USP
(2º andar – em frente ao elevador)

Horário de atendimento:
Das 13:00 às 19:00, de segunda a sexta

Conheça a gestão: 

\textbf{Professoras}: 
\vspace{-15pt}
\begin{itemize}
  \item Cristina Brech ({\tt brech@ime.usp.br})
  \item Renata Wassermann ({\tt renata@ime.usp.br})
\end{itemize}

Para mais informações, envie um e-mail para {\tt cip@ime.usp.br} ou acesse:
\begin{itemize}
  \item Site: \url{https:www.ime.usp.br/cip}
\end{itemize}


\end{subsecao}

% CinIME -----------------------------------------------------------------------
\filbreak
\begin{subsecao}{CinIME}

\figurapequenainline{novo-logo-cinime}

Convidamos vocês a conhecer o CinIME!! Somos um projeto do CAMat que promove a
exibição de filmes para discentes, docentes e funcionários do IME. Queremos pautar
de forma mais crítica o audiovisual, oferecendo um momento reflexivo, mas que 
também seja de lazer e diversão.

O CinIME ocorre toda sexta-feira. A sessão, o refrigerante e a pipoca são de graça. 
A organização é feita a partir de uma comissão aberta ao público, que
qualquer estudante pode compor e trabalhar conjuntamente na construção do projeto. 

Além das sessões usuais, também fazemos projetos paralelos, como a MostrIME - Mostra de
Animação no IME, em que exibimos dezenas de filmes de animação durante as férias, o Corujão
Satoshi Kon, que trouxe três filmes desse diretor exibidos numa madrugada, no CCSL (sim, a gente
passou a madrugada no IME vendo filmes!) e a nossa exibição do Oscar, em que escolhemos um filme
não-indicado para homenagearmos. 

A comissão organizadora do CinIME atualmente é guiada pela seguinte questão: "Quais 
horizontes imaginativos o cinema pode proporcionar aos estudantes do IME?". Mais detalhes
sobre isso podem ser encontrados no nosso projeto, o documento mais importante sobre o 
CinIME, que está disponível no site do CAMat.

Não se esqueçam de sugerir filmes, votar e comparecer ao CinIME!

Sigam as redes do CAMat, entrem na nossa comunidade no Discord e fiquem por dentro das novidades!

\begin{description}
  \item[Discord:] \url{https://discord.gg/qDfXUMVm6j}
  \item[Site:] \url{https://camat.ime.usp.br/cinime/}
  \item[Instagram:] \url{https://www.instagram.com/camat.usp/}
\end{description}

\end{subsecao}


% Grupo A5 ---------------------------------------------------------------------
\begin{subsecao}{Grupo S4}

\begin{wrapfigure}{r}{0.25\textwidth}
    \vspace{-25pt}
    \begin{center}
      \includegraphics[width=0.25\textwidth]{img/logo S4 maior.png}
    \end{center}
    \vspace{-25pt}
  \end{wrapfigure}


Nós somos um grupo de estudantes de diferentes anos da graduação e pós-graduação 
do IME, e organizamos seminários feitos por alunes e para alunes em diversos
tópicos de matemática, muitos não vistos na graduação. Originado em 2015 na forma
de seminários de iniciação científica, o S4 tem hoje em dia o propósito de apresentar seminários 
periódicos acessíveis e interessantes para estudantes que ainda não possuem muita 
bagagem matemática, inclusive de outros cursos da USP. 

Somos independentes de outras instituições do IME, e estamos sempre abertos a 
receber novos participantes, dos mais engajados aos mais casuais. É um espaço
para que es alunes tenham a oportunidade de apresentar algum tema que tenham
estudado, e começar na prática a treinar a habilidade de comunicação científica. 
Es alunes, independentes do ano, estão convidades a sugerir temas e dar ideias
para a organização de seminários. Caso desejem apresentar um seminário, nós
poderemos ajudar no seu preparo em um ambiente amistoso e receptivo.

Além disso, o S4 também é um ambiente de estudo e socialização, onde todes podem
compartilhar dúvidas, sugestões, conhecer novos bixes e veteranes, bater papo,
descontrair, etc. Estamos sempre dispostos a ajudar quaisquer alunes, especialmente
bixes, para que possamos ver um pouco mais do que a matemática tem a oferecer e
compartilhar as diferentes visões sobre ela de cada um de nós.

\begin{description}
  \item[Discord:] \url{discord.gg/XzXeCzc3rk}
  \item[E-mail:] \url{grupo-s4@riseup.net}
\end{description}

\end{subsecao}


% Olimpíadas de Matemática e Informática ---------------------------------------
% Maratona de Programação ------------------------------------------------------
\begin{subsecao}{Olimpíadas de Conhecimento}

\begin{itemize}

\item{\bf Matemática: }

Bom, pessoal, se vocês entraram no IME, podem ter ouvido falar ou participado
de alguma Olimpíada de Matemática no Ensino Fundamental e/ou Médio. A
boa notícia é que vocês vão poder continuar participando se quiserem,
e quem nunca participou tem a oportunidade de começar agora.

Mas por que participar? As Olimpíadas Universitárias de Matemática são uma
oportunidade de se divertir resolvendo problemas difíceis de Matemática e agregar
valor ao currículo ao mesmo tempo. Elas são parecidas com as Olimpíadas de
Ensino Médio, mas com conteúdo de Matemática da graduação (essencialmente
Cálculo, Análise, Álgebra Linear, Álgebra, Combinatória e Teoria dos Números),
mas com enfoque em problemas que exigem criatividade e técnicas mais inovadoras,
muitas das quais vocês provavelmente não verão durante toda a graduação.

De quais olimpíadas podemos participar? Como alunos de graduação, vocês podem
participar da Olimpíada Iberoamericana de Matemática Universitária (OIMU),
Olimpíada Brasileira de Matemática (OBM) e Olimpíada Internacional de
Matemática (IMC).

Como fazemos para nos preparar? Os sites institucionais dessas olimpíadas
têm material voltado para vocês que querem estudar e se preparar, mas vocês
podem procurar alguém mais experiente para indicar alguma bibliografia por fora
também. 

Como fazemos para participar? Inscrevam-se pelo site ou entrem em contato com
o professor Colucci. Para o IMC aconselha-se ter ganhado medalha na OBM,
já que é necessário apoio financeiro do IME por ser uma olimpíada internacional.

%REFTIME
Mas nós, bixes, temos chance? Como foi o desempenho de IMEanos nelas? A organização 
da OBM criou a Copa Elon Lages Lima, que é a primeira fase da OBM, mas tem uma 
premiação própria e um nível mais acessível para quem chegou agora. Desde 2021, 
vários alunos do IME conseguiram medalhas, inclusive bixes que foram chamados para participar 
da OBM! É a melhor porta de entrada para quem ainda não tem muita experiência. 

Se tiverem alguma dúvida, não hesitem em perguntar a algum veterane sobre os
Olímpicos!

Links institucionais:

\begin{description}
  \item[] \url{http://oimu.eventos.cimat.mx}
  \item[] \url{https://www.imc-math.org.uk/}
  \item[] \url{http://www.obm.org.br}
\end{description}

\item{\bf Informática: }

\textit{``Informática? Vocês mexem com Word, Excel e PowerPoint então?''}

Responder essa pergunta já virou rotina para competidores da Olimpíada
Brasileira de Informática (OBI). Não, Informática não é Word. Oras, então o que é a OBI?

A OBI é uma competição de lógica, matemática e computação. As provas envolvem
alguns problemas que você deve resolver com programas de computador.

Esta competição, na graduação, é exclusiva para ingressantes recém formados do
ensino médio. Quer dizer que vocês são a nossa única esperança de trazer mais
gloriosas medalhas ao IME! Isso também quer dizer que essa é a sua única chance
de participar da OBI, uma competição relativamente tranquila comparada à
Maratona de Programação.

Para participar, basta falar com o MaratonUSP
(\url{https://www.ime.usp.br/~maratona/}), um grupo de extensão focado nesse
tipo de competição que promete te ajudar a se inscrever e se preparar, ou com o
Professor Carlinhos (\url{http://www.ime.usp.br/~cef/}).

%\item{\bf Maratona de Programação: }

%À primeira vista, a Maratona de Programação pode soar um tanto
%surreal. Nerds correndo pela USP ao mesmo tempo que resolvem
%problemas de computação e matemática? Infelizmente esse não
%é o caso.

%A Maratona de Programação se resume à resolução de problemas.
%Se você adora resolver desafios, quebrar a cabeça com novos
%e excitantes problemas e acumular toneladas de dinheiro, esse
%é o lugar perfeito para você!

%A competição consiste em uma série de problemas que englobam
%temas como programação dinâmica, grafos e estruturas de dados.
%Times de três pessoas devem resolver a maior quantidade de
%desafios em cinco horas de programação. E tudo isso com direito
%a um lanche gratuito durante a prova.

%Mas não temam, bixes. Não é só por que vocês acabaram de entrar que
%a probabilidade de ganhar uma medalha seja nula. Inclusive, na primeira
%fase da maratona, uma equipe de bixes tem vaga garantida para a
%fase brasileira.

%Além da fama, constantes pedidos por autógrafos e dinheiro de sobra,
%a maratona também vai lhes trazer um conhecimento muito mais
%adiantado em relação ao dos seus colegas de classe, e até oportunidades
%de emprego em empresas de renome, como Google, Facebook e IBM.

%Se vocês se interessaram pela maratona e querem saber os horários dos
%treinos, como participar ou saber mais, acessem:

%\begin{description}
%  \item[Site:] \url{http://www.ime.usp.br/~maratona}
%  \item[Site da competição nacional:] \url{http://maratona.ime.usp.br/}
%\end{description}

\end{itemize}


\end{subsecao}


% Fala Sério -------------------------------------------------------------------
%\input{fala_serio.tex}

% Bee Data USP -----------------------------------------------------------------
\input{beedata.tex}

% Coloquinho da Pós-Graduação ---------------------------------------------------
\begin{subsecao}{Coloquinho da Pós-Graduação}

\figurapequenainline{coloquinho}

O Coloquinho da Pós-Graduação foi criado para proporcionar um espaço de discussão 
acadêmica e troca de conhecimentos entre os pós-graduandos e o público geral. 
As apresentações orais abrangem uma ampla variedade de tópicos relevantes da 
Matemática Pura e Aplicada. 

No 1° semestre de 2025, os seminários serão às \textbf{terças-feiras} às
\textbf{13:00}, no IME-USP. O primeiro seminário deste ano ocorrerá no 
dia 19/03/2025. 

Todes são bem-vindes para assistir os seminários. Emitimos certificados 
de participação para as horas complementares. 

Acessem o nosso site e nos sigam no Instagram para verem as apresentações
anteriores e acompanharem as atualizações sobre as datas e o local dos próximos 
seminários!

\textbf{Instagram:} @coloquinho.ime
\textbf{Email:} @coloquinhoime@gmail.com
\textbf{Site:} https://sites.google.com/ime.usp.br/coloquinho 



\end{secao}


% CCSL - Centro de Competência em Software Livre

\pagebreak
\input{ccsl.tex}

% CRInt ------------------------------------------------------------------------
\input{crint.tex}

% Cursos de Verão --------------------------------------------------------------
\begin{secao}{Cursos de Verão}

o IME, entre Janeiro e Fevereiro, são oferecidos diversos cursos de verão,
dentre os quais estão presentes cursos de difusão cultural e disciplinas de
pós-graduação. Os cursos são abertos não apenas para a comunidade IME, mas
também para a comunidade externa.

Você pode fazer cursos de tecnologia caso queira aprender uma linguagem ou ferramenta nova, 
ou então acrescentar uma formação complementar no currículo. Também existem cursos mais teóricos 
que podem te ajudar com as disciplinas do curso, ou então te apresentar uma área nova da matemática. 
Tem até cursos como o de resolução de problemas e criatividade, que você pode fazer só por diversão. 
Para se inscrever é cobrada uma taxa, mas se você precisar pode pedir uma isenção. 

Se vocês estiverem muito empolgados durante as férias ou sem ter o que fazer, fiquem atentos também para 
cursos de verão em outras unidades da USP, como o IF, o IAG e a FFLCH (você pode, por exemplo, resolver aprender 
japonês durante as férias)! Fora da USP, os cursos do IMPA também atraem muitos IMEanos nessa época do ano, tanto 
pela qualidade quanto por aproximar os estudantes de uma pós no IMPA. Como o IMPA fica no Rio de Janeiro, é preciso 
alugar uma casa para passar o período do curso lá, mas também é possível solicitar uma ajuda financeira caso vocẽ precise.


Saiba mais em: 
\url{https://www.ime.usp.br/verao/}
\url{https://www.ime.usp.br/pos/verao/}
\url{https://impa.br/ensino/programas-de-formacao/programa-de-verao/}
\url{https://sce.fflch.usp.br/node/3281}

\end{secao}


% CEC --------------------------------------------------------------------------
\begin{secao}{CEC - Centro de Ensino de Computação}

\begin{subsecao}{O que é?}

O Centro de Ensino de Computação (CEC) foi criado em 1993, pelo
Departamento de Ciência de Computação do IME, com o objetivo de
oferecer apoio às aulas dos cursos de graduação do Instituto e para os
treinamentos/cursos oferecidos à comunidade USP e não USP, através de
seus laboratórios com recursos de informática.
Está localizado no Bloco B do IME, ao lado da Seção de Alunos. O horário
de funcionamento do CEC é de 2a a 6a das 8h00 às 22h00 e nos recessos
escolares o horário de abertura é alterado para 9h00.

\end{subsecao}

\begin{subsecao}{Estrutura}

O CEC possui 137 computadores, distribuídos em 04 (quatro) salas/laboratórios,
sendo 02 (duas) delas com capacidade para até 60 pessoas e com prioridade de uso
para as aulas de graduação do IME e para os cursos/treinamentos. Os outros
laboratórios são utilizados por alunes de graduação dos cursos do IME para a
realização dos trabalhos e pesquisas acadêmicas. Em todos os computadores o Sistema
Operacional é o Linux Debian e possuem os aplicativos solicitados pelos docentes
do Instituto.

\end{subsecao}

\begin{subsecao}{Como funciona?}

Possui rede cabeada e para utilizá-la é necessário ter uma senha de
acesso (conta de usuário). Você precisa ser alune matriculado no IME para
solicitar a senha de uso, pois o CEC não é uma sala pró-aluno (para ver como as 
salas pró-aluno funcionam acesse \url{https://www5.usp.br/keywords-s/salas-pro-aluno/}). 
Para solicitar a senha de acesso, envie e-mail para \url{cec-senha@ime.usp.br},
com o assunto "CEC senha". Informe seu nome completo e número USP. Utilize
seu e-mail institucional (@usp.br - se esqueceu entre em
\url{https://id.usp.br/}).

\end{subsecao}

\textbf{Obs.:} bixes, ao frequentarem o CEC, fiquem atentes ao ar-condicionado. Se
estiver ligado, deem preferência a usarem calças, blusas, jaquetas e meias de
lã. Cobertores são opcionais. Se não, boa sorte ou \textit{hasta la vista}!

\end{secao}

\pagebreak
% CAEM -------------------------------------------------------------------------
\input{caem.tex}
\pagebreak

% Um Pouco Sobre a USP ---------------------------------------------------------
\begin{secao}{Um Pouco Sobre a USP}

Vocês, que estão conhecendo a USP pela primeira vez, devem saber desde cedo que aqui
há muita burocracia. É bom que se acostumem com ela, já que vocês terão que enfrentá-la.

O reitor é presidente do maior órgão da USP, o Conselho Universitário (abrevia-se
C.O. para evitar frases do tipo: ``Vou ter uma reunião no CU hoje'', ``O CU não
está funcionando muito bem esse semestre'', ``Os alunos não têm acesso ao CU''
etc.) que determina TODOS os rumos da universidade.

Abaixo dele vêm as coordenadorias e unidades. O SAS (Superintendência de Assistência Social),
por exemplo, é o departamento responsável pelos serviços
que a universidade oferece (não é a melhor coisa do mundo, mas oferece) para a
comunidade universitária: ônibus circulares, alimentação através dos restaurantes
universitários (famosos bandejões), moradia para estudantes (CRUSP) etc.

Alguns outros lugares que vocês devem saber que existem são o HU (Hospital
Universitário), o CEPE (Centro de Práticas Esportivas, leia ``cepê''), o banheiro
da FEA (Faculdade de Economia, Administração e Contabilidade), que fica na
frente do IME e é uma das faculdades mais bem abastecidas financeiramente
(apelidada ``carinhosamente'' de Shopping), a Física (certas aulas de laboratório são lá)
 e a querida Faculdade de Educação (para o pessoal da Licenciatura).

\begin{subsecao}{e-Disciplinas e E-mail USP}

Bixes, atenção e cuidado com o e-mail que vocês criaram (ou vão criar)! É um
email@usp.br, do sistema ``oficial'' de e-mails da USP, o que significa que é
através dele que a Universidade, o Jupiterweb, a diretoria do IME, o CAMat e
alguns desocupados lhes enviarão comunicados oficiais, o que pode envolver as
oportunidades diversas de intercâmbios, estágios, monitorias e todas as outras coisas
que vocês, simples bixes IMEanos, sejam capazes de se imaginar fazendo. Por este e-mail
vocês receberão informações antes mesmo das principais serem colocadas nos paineis eletrônicos que
ficam espalhados pelo IME (é, nem todas são colocadas nesses paineis não).
Com as aulas on-line ou não, o e-mail sempre foi uma das principais formas de
comunicação (às vezes a única) entre nós alunos e todos os outros grupos, pessoas,
órgãos etc. que compõem as comunidades IMEana e da USP, então é sempre bom verificar
a caixa de entrada do seu e-mail com frequência! 

Além disso, peçam para os veteranes falarem mais um pouco de algumas 
(das \textbf{muitas}) vantagens de se ter um e-mail USP e aproveitem tudo aquilo que 
vocês agora têm direito!

A USP tem, há alguns anos, um ambiente virtual de aprendizagem para dar apoio
às disciplinas: o Moodle da nossa universidade, oficialmente conhecido como
e-Disciplinas (é comum chamarmos dos dois jeitos e há alguns professores que
se referem a ele como ``Stoa'', que é um nome bem antigo!). Você pode acessá-lo em
\url{https://edisciplinas.usp.br/}, fazendo login com o e-mail USP e a senha única.

Na página inicial, você provavelmente encontrará links para as páginas da maioria
das disciplinas que vai cursar (se não todas). Muitos professores não utilizam o 
e-Disciplinas como apoio aos seus cursos, às vezes fazem isso nos seus próprios 
sites institucionais ou simplesmente não usam a internet para nada. 

\end{subsecao}

\begin{subsecao}{Grupos no Whatsapp}

Com o intuito de organizar as atividades referentes a uma disciplina e ser uma
maneira prática e efetiva de passar recados, costumamos utilizar grupos do WhatsApp
para cada uma das matérias. Caso os links ainda não tenham sido criados,
vale a pena juntar seus amigos para criar os grupos no começo do semestre e espalhar
para os outros alunos.

Dentre as principais funções do grupo do WhatsApp estão:

\begin{enumerate}
%\item Colocar o link para acessar a plataforma on-line das aulas na descrição pra facilitar pra todo mundo;
\item Compartilhar os materiais e listas disponibilizados pelo docente da disciplina;
\item Passar recados importantes, sejam datas para entrega de trabalho, provas etc.;
\item Tirar dúvidas sobre a matéria com os coleguinhas.
\end{enumerate}

Pra ficar por dentro de tudo o que acontece na matéria, não se esqueça de procurar o
link e fazer parte dos grupos de whats das matérias que você está cursando!! :)

\end{subsecao}

\begin{subsecao}{Eduroam}

O Eduroam (EDUcation ROAMing, não tem nada a ver com alguém chamado Eduardo) é um
serviço de wi-fi que \sout{às vezes} provê acesso à Internet dentro de diversas
universidades (inclusive a USP). Mesmo que você seja apenas aluno da USP, se viajar
para outra universidade, no Brasil ou no exterior, que possui o Eduroam, ainda poderá
ter acesso à Internet usando o mesmo login!

Para usar o Eduroam, basta se conectar à rede \texttt{eduroam}, colocando no campo
de login o seu número USP seguido por @usp.br. Por exemplo, se o seu número USP for
\texttt{123456789}, então deverá colocar no campo de login \texttt{123456789@usp.br}.
No campo da senha, deverá colocar a sua senha única (a mesma senha que vocês usam
para logar no Jupiterweb).

Tem várias outras informações sobre as configurações na hora de fazer login que
\sout{ninguém entende, mas se algum de vocês entender,} estão no site
\url{https://eduroam.usp.br/como-usar/}. Nesse site também tem alguns programas para Windows que 
configuram o Eduroam automaticamente. (Dica: se for baixar eles, utilize o Instalador Alternativo) %REFTIME
Também existem cartazes espalhados pelo IME com os detalhes de login para diferentes tipos de 
dispositivo, então deem uma olhada nas paredes.

\end{subsecao}

\begin{subsecao}{O SAS}

A Superintendência de Assistência Social (SAS), que fica próxima à Praça do
Relógio, é o órgão da USP responsável pelo bem estar financeiro dos alunos — não
é lá a melhor coisa do mundo, mas ajuda... É onde vocês podem conseguir seus
benefícios, tais como auxílios financeiros, moradia gratuita (CRUSP)
ou auxílio alimentação (conhecido como vale-bandex) e dentistas gratuitos.
A SAS também é responsável pelo Setor de Passes Escolares.

Se algum de vocês, bixes, estudou a vida inteira em escola pública, tem baixa renda,
ou está passando por algum tipo de vulnerabilidade social, não pense que está sozinho nessa,
muitos também passaram por isso e vocês conseguem encontrar apoio aqui. Permanência
estudantil é um direito e não deve ser motivo de vergonha, portanto busquem
auxílio, pois podemos ajudar. Seguem as diversas alternativas para todos 
que precisam;

\textbf{Moradia e auxílios financeiros}

Se vocês infelizmente não têm tanto acesso a meios de locomoção, ou dinheiro para
pagar transportes ou mesmo repúblicas, saibam que diversos auxílios podem ser
oferecidos para vocês para amenizar a sua situação:

% Com a carteirinha (e-card) que você recebeu na matrícula, você deverá entrar
% na página do SAS (\url{http://sites.usp.br/sas/}) e solicitar a inscrição
% para o processo de moradia e alojamento. Lá você terá um formulário, três
% trilhões de coisas para preencher e assim que te chamarem, deve apresentar os
% documentos (hum!) necessários para a assistente (você irá “ganhar” uma). Na
% improvável hipótese do site estar fora do ar, como todo ano acontece, você terá
% que ir ao SAS e pedir alojamento na USP, lá no Bloco G do CRUSP.

%REFTIME

As inscrições dos ingressantes para o Programa de Apoio à Permanência e Formação Estudantil (PAPFE)
deverão ser feitas dentro dos prazos estipulados no EDITAL PAPFE 2025, que pode 
ser encontrado na página do SAS (\url{https://prip.usp.br/wp-content/uploads/sites/1128/2024/12/EDITAL-PAPFE25-v-Final.pdf}).
No Portal de Serviços Computacionais da USP \url{https://portalservicos.usp.br/}, vocês entrarão 
com os dados da sua conta USP para se inscrever, responderão a questionários e anexarão (um bilhão de) 
documentos comprobatórios da situação socioeconômica declarada no ato de inscrição.

Vale ressaltar, bixes, que vocês tomem o cuidado de não fazer o mesmo que alguns
bixes de outros anos, que confundiram alojamento (vaga no CRUSP) com o auxílio moradia (financeiro). São dois
requerimentos distintos e é bom que vocês fiquem atentos na hora de se inscrever, verificando
qual(is) tipo(s) de auxílio(s) vocês precisam.

Os requerimentos conterão perguntas sobre a situação socioeconômica, bem como os
documentos que deverão ser trazidos para comprová-la. Perguntas como renda
salarial, quantas pessoas contribuem com ela, números de bens móveis e imóveis,
situação habitacional, tipo de escola em que estudaram, se trabalham e há quanto
tempo, quanto gastam para vir à USP, o tempo de ida etc. Ainda há um espaço
para descreverem alguma particularidade não exposta nas perguntas que, obviamente,
receberá um parecer técnico.

No ato da inscrição, você pode optar entre as ruas modalidades mencionadas: o auxílio permanência
no valor de 850 reais, ou a vaga no alojamento (CRUSP), acompanhada de um auxílio de 320 reais. Como 
o número de vagas no CRUSP é pequeno, naturalmente essa modalidade é mais concorrida que o auxílio financeiro integral.
Caso você seja selecionado em qualquer uma dessas modalidades, recebe também o auxílio alimentação, que se trata de créditos
para carregar o cartão para o bandejão gratuitamente, todo mês, através do JúpiterWeb.

%\textbf{Bolsa-trabalho}

%Destina-se a alunos de graduação vinculados a projetos de extensão de serviços à
%coletividade. Os projetos são selecionados anualmente, de acordo com sua relevância
%para as finalidades da universidade pública e os estudantes vinculam-se por
%afinidade acadêmica ou científica. Cada bolsa é de 1 (um) salário mínimo por
%40 horas de trabalho mensais. Além da seleção socioeconômica feita pela
%DPS (Divisão de Promoção Social), há uma seleção técnica feita pelos supervisores
%dos projetos.

%Mas fiquem espertos! Para tudo tem prazo e o SAS não é obrigado a ficar esperando a
%boa vontade de aparecerem por lá de ninguém. Qualquer dúvida, liguem pro SAS.

\textbf{Atendimento odontológico gratuito}

Antigamente para vocês agendarem o atendimento odontológico gratuito era necessário fazer uma
carteirinha no HU (Hospital Universitário), em seguida comparecer ao Bloco G
do CRUSP com a carteirinha e agendar, porém hoje em dia, graças à tecnologia,
isso não é mais necessário: você só vai precisar mandar uns e-mails, preencher uns
formulários e esperar \sout{uma eternidade}.

Primeiro você precisa mandar um e-mail solicitando triagem para {\tt triagemodonto.sas@usp.br},
vão te responder com um formulário que você vai preencher e enviar de volta, após isso eles
te colocarão na fila de espera pra triagem. No final, deve demorar muitos meses pra
eles te chamarem e quando fizerem isso você só vai ter que comparecer ao Bloco G do CRUSP
munido do seu cartão USP.

No final será demorado, porém o serviço é gratuito e realmente ajuda. Eles também atendem casos
de emergência como por exemplo dor de dente, algum dente quebrado ou qualquer outra situação
do gênero. No caso de dúvidas, são poucos os que podem te ajudar nesse assunto, então o melhor
a se fazer é ligar pra lá: (11) 3091-3393.

%FIXME
\textbf{Setor de passe escolar}

Será nesse lugar que você resolverá boa parte dos seus problemas com o SAS e com o
Bilhete Único. Você pode se dirigir ao Bloco “E” do CRUSP para receber orientação sobre
as mais diversas burocracias em que está se metendo.

Para as linhas da EMTU, SPTRANS, METRO, vocês devem entrar no site do SAS e fazer
o pré-cadastro para o respectivo cartão (\url {http://sites.usp.br/sas/}) para que a USP envie os 
seus dados para as empresas de transporte. Talvez demore um pouco, pois a USP tem que
avisar para a SPTrans que vocês passaram na Fuvest. Fiquem atentos! Qualquer dúvida,
liguem para a seção de passe escolar do SAS: (11) 3091-3581. Vocês também podem entrar nos sites 
das empresas de transporte e procurar por mais informações lá.

Há também um cartão especial para alunos de universidades públicas provenientes
de outras cidades do interior de São Paulo. Se algum de vocês vem de algum desses domos
ignotos (como Resende, Caçapava ou Guaíra), pode se dirigir ao guichê da sua empresa de
transporte intermunicipal (Cometa, Danúbio Azul etc.), apresentar seu Cartão USP,
preencher um formulário e eles aguardarão confirmação do seu instituto. Então
você pega seu cartão na Seção de Alunos e, na compra de passagens entre São Paulo
e sua cidade-natal, paga 50\% do valor normal. Só não deixe isto por último na
sua lista de necessidades porque existe um período do ano em que os guichês
liberam seus formulários; em resumo, espiche suas orelhas e corra para a rodoviária.

\end{subsecao}

\begin{subsecao}{Atendimento Psicológico}

Existem inúmeros serviços na USP que oferecem atendimento psicológico
aos alunos, cada um com um foco e um público-alvo. Dentre eles, destaca-se
o Serviço ECOS, um serviço focado em três etapas: Escuta, Acolhimento e Ação.

O ECOS conta com um espaço de atendimento situado no campus Butantã, realizando 
escutas pontuais presenciais e sem necessidade de agendamento para reflexão e possíveis 
encaminhamentos às redes de cuidado interna ou externa à USP. As escutas são realizadas 
pela equipe ECOS, composta por profissionais de saúde e alunes bolsistas de diversas áreas 
do conhecimento, que desenvolvem projetos de iniciação científica. O programa não atende
emergências. Para situações de crise aguda em saúde mental, deve-se buscar um Pronto Socorro 
Psiquiátrico (a referência para o campus Butantã é o Pronto Socorro Lapa, Av. Queiroz Filho, 
313; telefone: 4878-1701).

O Programa articula continuamente uma rede de cuidados em saúde mental baseada nos
territórios dos campi USP, mobilizando ações e serviços universitários, do SUS, de 
outras políticas públicas e da sociedade, que visa garantir acesso a estes recursos 
pela comunidade universitária e fornecer um Mapa da Saúde Mental da USP.

A equipe do ECOS oferece apoio institucional a diferentes grupos, unidades e campi da 
USP, seja para a criação de ações e serviços relacionados à saúde mental, seja para a 
qualificação de iniciativas já existentes. Dessa forma, realiza o acolhimento a demandas 
coletivas, buscando potencializar ações locais e a participação das respectivas comunidades 
na promoção da saúde mental em cada território.

Caso você esteja precisando, o local é próximo ao bandejão Central, confira o endereço, contato e informações: 
Endereço: Rua do Anfiteatro, 181 – Favo 22 – Cidade Universitária – São Paulo/SP (em frente à portaria do Bloco C do CRUSP)
Horário de funcionamento: Segunda a sexta-feira, das 08h às 17h
Telefone: (11) 3091-8345
E-mail: ecos.prip@usp.br

Uma outra opção são os serviços oferecidos pelo Instituto de Psicologia (IPUSP).
Os principais serviços do IP são o LEFE, o Atendimento Presencial e o PAPO:

O LEFE se trata de um plantão psicológico oferecido à comunidade, em que
você pode conversar com um profissional num momento de crise. O plantão 
funciona na terça-feira às 17 h, no IP, e para ser atendido basta chegar 
entre as 16 e 17 h.

Mais informações sobre o LEFE: 
\url{https://www.ip.usp.br/site/plantao-psicologico-lefe/}

O Atendiento Presencial no IP é o serviço em que você pode ter acesso a
atendimento individual de forma contínua por um período mais longo. Essa
modalidade também exige inscrição, e os inscritos ficam em uma lista para
passar por triagem. O período com maior chance de ser selecionado é no inicio
do semestre, mas também é possível conseguir vagas em outros momentos em 
situações excepcionais. 

Mais informações sobre o Atendimento:
\url{https://www.ip.usp.br/site/atendimento-4/}

O PAPO é o Projeto de Apoio Psicológico Online, que foi criado durante a 
pandemia para oferecer de forma segura uma via de atendimento psicológico
para a comunidade. Para essa modalidade de atendimento é preciso se inscrever
previamente, no período indicado no site.
\url{https://www.ip.usp.br/site/covid-19-apoio-psicologico-online/}

O IP também oferece outros serviços, como o Ateliê Aberto e o Serviço
de Orientação Profissional. Você pode conhecer eles aqui:
\url{https://www.ip.usp.br/site/centro-escola-do-instituto-de-psicologia-da-usp-ceip/}


Bixes, sempre que precisarem, não tenham receio de procurar apoio
psicológico, por mais banal que vocês possam achar o motivo. Vocês
podem frequentar os serviços oferecidos pela USP para conversar sobre
qualquer coisa que esteja lhes incomodando, sejam professores,
matérias, EPs ou algo do tipo sem criar qualquer
problema para vocês no IME.

\end{subsecao}

% Utilize para começar uma nova página do lado esquerdo do Guia!

\begin{subsecao}{Bandejão}
%TODO Fazer as duas páginas ficarem na mesma "visualização" no guia impresso
%Ou seja, usar esquema Odd/Even do LaTeX.
% Isso foi deixado de lado nessa versão de 2016 para que o mapa dos circulares
% ficasse no lugar certo e não tivesse nenhuma página em branco
% (como o cleardoublepage ali em cima)
% 2017: não sei se entendi o que era para ser feito, mas usei o
%       \cleardoublepage para a imagem do calvin não ficar zuada.


Os bandejões, vulgarmente conhecidos como Restaurantes do SAS, são os lugares
em que vocês podem se alimentar razoavelmente a um preço analogamente razoável.
Os tickets custam R\$2,00, sendo estes carregados na carteirinha USP. Antigamente, o
lugar para carregar o ``bandejão único'' era no CARE (Centro USP de Acolhimento e 
Referência para Estudantes), perto do SAS e do bandejão Central, porém, em 2024, essa
opção deixou de funcionar. Além disso, é possível verificar o saldo atual e comprar 
mais créditos via Pix no site {\tt https://uspdigital.usp.br/rucard/} ou pelo
aplicativo ``Cardápio+ USP'', disponível para Android e iOS (não se esqueça de olhar 
também os outros aplicativos da USP disponíveis, pois podem ser muito úteis).
Segundo os meios oficiais (e-mails enviados pela usp), a opção de boleto não 
estaria mais funcionando, porém, alunos relataram que essa opção ainda aparece 
para eles. Portanto, caso queiram e consigam recarregar o RUcard via boleto, fique 
sabendo que os créditos demoram até 03 (três) dias úteis para estarem disponíveis.
Se realizarem o pagamento pelo Pix, os créditos podem ser usados dentro de alguns
minutos.
Dica: Confira o saldo da carteirinha no site ou no aplicativo antes de tentar usar!

Existe uma lenda que sobremesas especiais aparecem em certos bandejões perto de
datas comemorativas. Mas não se engane - normalmente, elas não aparecem no
cardápio (talvez para não ficar muito cheio de gente).

\figuragrande{bandex_calvin}

O cardápio semanal do bandejão pode ser visto no site {\tt
http://sites.usp.br/sas/} ou pelo aplicativo ``Cardápio+ USP'', mas se vocês
estiverem com preguiça de ver em um desses lugares, é bem provável que um veterane 
já saiba e resolva informar, se questionado com muita educação.

O cardápio é geralmente composto de arroz (com opção normal e integral), feijão,
prato principal (carne/ovos), acompanhamento (legumes ou verduras refogadas,
cremes, molhos), salada (algumas folhas), sobremesa (frutas e/ou, às vezes, algo mais
requintado), pãozinho (normal e integral) e suco (de diversos sabores: azul, vermelho,
amarelo, etc.), além de temperos genéricos (jamais perguntem do que eles são feitos). 
Se vocês são vegetarianos, o acompanhamento nunca contém carne (e estão tentando 
fazer com que sejam veganos também), e todos os bandejões têm uma alternativa 
vegetariana de PVT para o prato principal, normalmente vindo na forma de ração,
podendo também aparecer nas formas de escondidinho, lasanha, kibe e outros
(vale a pena experimentar o escondidinho da Central).


Não se esqueçam de levar as suas canecas do Kit-bixe se forem comer nos
bandejões, tanto para ostentar a posição de bixes do IME, quanto para salvar o
meio ambiente, evitando assim o uso de copos descartáveis. Os bandejões não 
disponibilizam copos e não aceitam o uso de garrafas nas máquinas de suco, por 
isso, é muito importante que vocês sempre tenham um copo ou caneca para conseguir 
tomar o líquido colorido saborizado do bandeco.


PS: O Efeito Bandex é proporcional à quantidade de salitre utilizado em cada
bandejão.\\
PS 2: Nunca, em hipótese alguma, jamais, visite a cozinha do seu bandejão de
preferência, pois você corre o risco de nunca mais almoçar na vida. Como já
dizia o velho sábio ``A ignorância é uma virtude''.\\
PS 3: Playstation 3.

Consultem a tabela abaixo para decidir em qual dos bandejões vocês vão comer.

\begin{figure}[!htbp]
\begin{center}
 	\includegraphics{img/tabela_bandeco_2023.pdf}
\end{center}
\end{figure}


O horário de funcionamento dos bandejões é o seguinte:\\
\textbf{Café da manhã:} 7h às 8h30 (somente no Central, de segunda à sábado)\\
\textbf{Almoço:} 11h15 às 14h15 (em todos os bandejões de segunda à sexta, de sábado 
somente do Central e de domingo somente na Química)\\
\textbf{Jantar:} 17h30 às 19h45 (somente na Química, na Física e no Central de segunda à sexta)\\

\textbf{IMPORTANTE:} o funcionamento dos bandejões pode ser alterado ao decorrer do ano,
principalmente nos feriados, por isso recomendamos que vocês fiquem atentos ao 
aplicativo e às placas informativas presentes nos bandejões.

Para mais informações, visitem o site do SAS: \url{http://sites.usp.br/sas/} ou
utilizem o aplicativo ``Cardápio+ USP''.

\end{subsecao}

\begin{subsecao}{Outros lugares para comer na USP}

Caso vocês não queiram comer no bandejão, seja porque estão com medo do
prato do dia ou simplesmente porque querem comer algo diferente (e
possivelmente de sabor melhor), saibam que há alguns outros lugares na USP onde
vocês podem comer; por exemplo, os que estão na seguinte lista.

%REFTIME
\textbf{Importante:} as informações referentes a valores não estão atualizadas com
os preços vigentes em 2024. Interpretem onde estiver escrito ``custa R\$ $x$'' 
como ``custa R\$ $x$ ou um tanto a mais''.
% Acredito que nenhum dos preços esteja atualizado, então generalizei

\textbf{Lanchonete da Física}

Fica um pouco antes do bandejão da Física, nela vocês podem almoçar pagando R\$
63,00 o quilo, sempre tem uma razoável quantidade de saladas e pratos
quentes, além disso, tem churrasco com uma boa variedade de opções.
Além do almoço por quilo, são vendidos alguns salgados e lanches como sanduíches e
beirutes. E logo em frente à lanchonete, há um lugar que vende cookies por R\$
8,50 (3 unidades). Vocês vão sentir o cheiro MARAVILHOSO dos cookies quando forem bandejar.

\textbf{Restaurante do IPEN}

Fica na parte de cima da rua entre o IF e o Parque Esporte Para Todos, para
entrar lá vocês precisam da carteirinha USP. O quilo custa R\$49,90 para sócios da marinha 
e R\$54,99 para não sócios. Infelizmente não há uma variedade muito grande de comidas lá e o
restaurante só abre para pessoas de fora às 13h.

\textbf{Lanchonete da FAU}

Localizada dentro da FAU assim que você sobe a rampa. Lá além de lanches é
servido almoço tanto por quilo (custa R\$43,90) como prato feito. Tirando o por
quilo, qualquer outra coisa deve ser paga no caixa antes de retirar. % Desde o ínicio da pandemia, 
%a Lanchonete da FAU se encontra fechada, mas pode ser que em breve ela volte a funcionar.
% Acredito que voltou, mas não tenho certeza se ainda tem o almoço por quilo. 
Além disso, no terceiro andar da FAU tem a mesinha de doces que os alunos deixam lá
para vender. A variedade e quantidade de doces varia bastante com o dia e o
preço deles costuma ser no máximo R\$2,50.

\textbf{Restaurante da FEA}

Fica atrás da FEA, lá o quilo tem bastante coisa diferente e algumas até que
sofisticadas, porém o preço do quilo lá é R\$85,90, provavelmente é o mais caro
da USP. Geralmente é frequentado por professores e funcionários da USP. Além do
self-service, há algumas opções de prato feito que custam R\$36,00.

\textbf{Trailers de lanche da \sout{ECA} Química}

Ficam na calçada com a Avenida Lineu Prestes, em frente ao Instituto de Química. 
Nesses carrinhos vende pastel, sanduíches, churros, tapioca, salgados, açaís, yakissobas e, 
inclusive, pratos feitos. Os preços dos lanches variam entre R\$11,00 e R\$22,00 e os das tapiocas 
entre R\$5,00 e R\$17,00. As opções de sanduíches e a quantidade de sabores de tapioca 
são variados e todos em geral são muito bons. Para os pasteis, vocês podem escolher até 5 recheios, 
e o preço varia de acordo com a quantidade de ingredientes escolhidos. O prato feito costuma custar 
entre R\$18,00 e R\$25,00 e o yakissoba entre R\$25,00 e R\$35,00.

\textbf{Cachorro-quente da Reitoria}

Fica na rotatória da biblioteca Brasiliana (ou a rotatória depois do CEPE, se assim preferir),
lá o cachorro-quente pode ser feito tanto no pão de cachorro-quente como na baguete. Além disso,
há a opção de pôr catupiry e de pôr uma salsicha adicional. O preço do cachorro-quente varia
entre R\$12,00 e R\$16,00.

\textbf{Poke da Reitoria}

Tem um lugar que vende Poke (um prato havaiano que mistura peixe cru em cubos,
arroz, shoyu e frutas) ao lado do cachorro-quente. Esta iguaria custa
aproximadamente R\$23,00 e você pode escolher os ingredientes que vão no seu prato,
tipo Spoleto. Dá pra fazer um Poke vegano, com cogumelos em vez de peixe. Além de
Poke, também vende Temakis, Hot Rolls, e até Açaí. A qualidade é boa, mas as filas
podem ser demoradas.

\textbf{Pastel do IEE}

Esse lugar maravilhoso fica a apenas alguns passos do IME. Apenas às terças-feiras,
em um canto escondido do IEE (pergunte a um veterane como chegar!) fica o famoso pastel do IEE.
Há muitas opções de sabores, além de poder comprar um caldo de cana para acompanhar (também tem 
várias opções de sabores!).

% \textbf{Restaurante do Viveiro}

% Esse restaurante foi inaugurado no final de 2023, e também fica muito perto do IME, próximo ao ponto 
% de ônibus ao lado do IO. Fica aberto das 11h às 14h30, e serve [...]
% Fica aqui um começo de texto se alguém quiser adicionar no futuro. 

%FIXME GAMBIARRA
\pagebreak
\end{subsecao}

\end{secao}


% Um Pouco Sobre o DCE ---------------------------------------------------------
\input{sobre_o_dce.tex}
\quadrinhos{5}
\pagebreak

% Hospital Universitário -------------------------------------------------------
\begin{secao}{Hospital Universitário}
   \begin{quote}\emph{O HU USP é o hospital de ensino de excelência utilizado
pelos Cursos de Atenção à Saúde da USP.  O hospital privilegia as pesquisas
relacionadas aos problemas de saúde  mais comuns da população brasileira. O
atendimento é regionalizado para o bairro do Butantã, sempre com enfoque no
ensino e pesquisa''}- trecho obtido da página do HU (\url{http://www.hu.usp.br/})
\textit{alguns} anos atrás.
   \end{quote}

Isso significa que os estudantes de Medicina, Ciências Farmacêuticas,
Odontologia, Saúde Pública, da Escola de Enfermagem e do Instituto de
Psicologia, mantendo contato direto também com os Institutos de Ciências
Biomédicas, de Biologia, de Química, a Faculdade de Arquitetura e Urbanismo,
a Escola Politécnica e a Escola de Comunicações e Artes precisam de cobaias para
suas atividades/experiências. O \sout{USPital} HU é um santo lugar que recebe alguns 
fracos de espírito que bebem demais e ficam incapacitados de fazer qualquer 
atividade fisiológica. 

Qualquer pessoa pode fazer seu cadastro no Hospital Universitário. Para isso, é 
necessário comparecer pessoalmente ao hospital portando os seguintes documentos:

\begin{itemize}
   \item a carteirinha USP;
   \item um documento oficial com foto (RG, CNH...);
   \item o CPF;
   \item o CNS (Cartão Nacional de Saúde, do SUS).
\end{itemize}

Estudantes da USP podem realizar esse cadastro a partir de um formulário 
disponibilizado no site do HU: (\url{http://www.hu.usp.br/}). É nesse mesmo site 
que você poderá agendar consultas e exames, cuide da sua saúde bixe!

Assim vocês possuirão alguns privilégios no atendimento, em situações de
emergência vocês são atendidos rapidamente (algo entre 600 minutos, como
diz na senha de espera). %e não precisam soletrar o nome da mãe enquanto
%estiverem desmaiados.

\begin{description}
\item [Como chegar:] Pegar 8012-10 ou 8022-10 sentido cidade universitaria
\item [Endereço:] Av. Prof. Lineu Prestes, 2565 - Butantã, São Paulo - SP, 05508-000
\item [Telefone:] (11) 3091-9200
\end{description}

Mas nem tudo são flores, HU não vai te atender para tudo, para muitos casos 
(99,9% das vezes) você vai ter que usar os serviços de uma UBS (Unidade Básica de Saúde),
caso seja de fora da cidade de São Paulo, confira no mapa qual UBS atende onde você mora
(\url{https://www.google.com/maps/d/viewer?mid=1-cDHamkJnuwbxrOEZVK7MLb7PoF2n88_&ll=-23.68533875879342%2C-46.521722701763906&z=10})
leve com você documento com foto, comprovante de residencia e o cartãozinho do SUS.

Feito o registro, a primeira coisa que você ira fazer, sera agendar um clinico geral para fazer um exame dermatologico,
com o laudo em mãos, você ira leva-lo ao CEPE para assim poder usar as piscinas de lá!

\end{secao}

\pagebreak

% Tudo Que Vai Volta -----------------------------------------------------------
\begin{secao}{Tudo Que Vai Volta (até bixes)}

\begin{subsecao}{Ônibus}

Se vocês não têm como ir nem como voltar, temos algumas dicas:

\begin{enumerate}
\vspace{-15pt}
  \item Trabalhem muito para comprar um carro, trabalhem mais para pagar a
  gasolina, venham para a USP de carro e, obrigatoriamente, deem carona a um
  veterane;

  \item Peçam a uma pessoa amiga para trazê-los e buscá-los durante seus
  longos anos de IME-USP;

  \item Conheçam alguém que, por sorte, more perto da sua casa, estude na USP,
  tenha o mesmo horário que você, seja legal e tenha carro.
  Traduzindo: s-o-n-h-e;

  \item Estiquem o dedão e esperem, esperem, esperem, esperem... a boa vontade
  de alguém para dar carona;

  \item Mudem-se para uma casa perto da USP;

  \item Peguem uma bicicleta e pedalem.

\end{enumerate}
\vspace{-15pt}
OK, vocês decidiram ser bixes normais e pegar ônibus! Nesse caso, vocês
provavelmente estão em uma das seguintes opções:
\vspace{-15pt}
\begin{enumerate}
  \item Vocês vêm de metrô pra USP;
  \item Vocês são super sortudos e acharam um ônibus que vai de casa pra USP.
\end{enumerate}
\vspace{-15pt}
Se vocês vão usar ônibus para ir ou voltar da USP, vocês precisam conhecer
os quatro pontos de ônibus perto do IME: FAU I, Oceanografia, FAU II e FEA.
Até 2019, esses quatro pontos tinham placas com os nomes escritos, mas
agora foram reformados e não têm mais as placas. Mesmo assim, a localização
deles não poderia ser mais intuitiva, veja só:

Tanto o ponto FAU I quanto o ponto da Oceanografia se localizam na Rua do
Matão -- FAU I é o ponto que fica mais próximo da FAU, e o da Oceanografia é o
que fica do mesmo lado da calçada que o IO (pronuncia-se i-ó), Instituto de
Oceanografia; os pontos FAU II e FEA ficam na Av. Prof. Luciano Gualberto (que a
partir de agora será chamada de Rua dos Bancos, para todo e todo o sempre),
e, como vocês já devem estar imaginando, são na frente da FAU e da FEA,
respectivamente. São muitos nomes para lembrar? Fiquem tranquilos: mais
cedo ou mais tarde vocês vão conhecer esses lugares e não vão nem se preocupar
mais com os nomes.

Como vocês viram ali em cima, nós damos apelidos também intuitivos para as
principais avenidas do campus:
\begin{itemize}
	\item Av. Prof. Luciano Gualberto = Rua dos Bancos;
	\item Av. Prof. Lineu Prestes = Rua do HU;
	\item Av. Prof. Mello Moraes = Rua da Raia.
\end{itemize}

Aqui está a lista de ônibus que passam em cada um dos pontos. Para mais detalhes
sobre cada linha, vocês podem usar o site da SPTrans ou o Google Maps.

\begin{subsubsecao}{Linhas municipais}

{\bf Ponto FAU I}

\begin{center}
	\begin{tabular}{|c|c|c|}
      \hline
	  Letreiro & Cor & Interligações (em ordem)\\
	  \hline
	  8012-10 - Metrô Butantã* & Laranja & Metrô: L4\\
      \hline
	\end{tabular}
\end{center}

{\bf Ponto da Oceanografia}

\begin{center}
	\begin{tabular}{|c|c|c|}
      \hline
	  Letreiro & Cor & Interligações (em ordem)\\
	  \hline
	  8022-10 - Metrô Butantã* & Laranja & Metrô: L4\\
   	  8082-10 - Metrô Butantã* & Laranja & Metrô: L4\\
      \hline
	\end{tabular}
\end{center}

{\bf Ponto FAU II}

(Esses os levam pra fora da USP)
\begin{center}
	\begin{tabular}{|c|c|c|}
      \hline
	  Letreiro & Cor & Interligações (em ordem)\\
	  \hline
	  177H-10 - Metrô Santana & Azul & Metrô: L4, L2, L3 e L1\\
	  7181-10 - Term. Princ. Isabel & Laranja & CPTM: L9\\
	  7411-10 - Praça da Sé & Laranja & Metrô: L4, L2, L3 e L1\\
	  7725-10 - Rio Pequeno & Laranja & - \\
	  809U-10 - Metrô Barra Funda & Laranja & L2 \\
	  8022-10 - Term. USP P3* & Laranja & -\\
      \hline
	\end{tabular}
\end{center}

Esses dois ônibus passam no ponto, mas estão CHEGANDO na USP, preste 
atenção!

\begin{center}
	\begin{tabular}{|c|c|}
	  \hline
	  Letreiro & Cor\\
	  \hline
	  701U-10 - Butantã-USP & Azul\\
	  702U-10 - Butantã-USP & Laranja\\
	  \hline
	\end{tabular}
\end{center}

{\bf Ponto FEA}

(Esses os levam pra fora da USP)
\begin{center}
	\begin{tabular}{|c|c|c|}
      \hline
	  Letreiro & Cor & Interligações (em ordem)\\
	  \hline
	  701U-10 - Metrô Santana & Azul & Metrô: L4, L2, L3 e L1\\
	  702U-10 - Term. Pq. D. Pedro II & Laranja & Metrô: L4, L2 e L3\\
	  7725-10 - Terminal Lapa & Laranja & CPTM: L8\\
	  8012-10 - Term. USP P3* & Laranja & -\\
      \hline
	\end{tabular}
\end{center}

Novamente, esses quatro ônibus passam no ponto, mas estão CHEGANDO na USP!
\begin{center}
	\begin{tabular}{|c|c|}
	  \hline
	  Letreiro & Cor\\
	  \hline
	  177H-10 - Butantã-USP & Azul\\
	  7181-10 - Cidade Universitária & Laranja\\
	  7411-10 - Cidade Universitária & Laranja\\
	  809U-10 - Cidade Universitária & Laranja\\
	  \hline
	\end{tabular}
\end{center}

As linhas marcadas com um * são as linhas circulares da SPTrans.
Segue abaixo suas descrições!

\end{subsubsecao}

\begin{subsubsecao}{Circulares}

Também conhecido como ``circulenda'' ou ``secular'' (aos sábados e domingos,
``milenar''), é o meio de transporte mais barato dentro da USP. Foi criado para
os USPianos se locomoverem dentro do Campus, mas em muitas vezes é melhor andar
do que ficar esperando. Existem 3 itinerários distintos: 2 deles têm trajetos
aproximadamente reversos e que cobrem todo o Campus, e o outro é uma
``versão expressa" de um deles, que só passa pelas partes do Campus que ficam
entre os portões P1 e P2 (isso inclui o IME, como vocês devem ter visto na
tabela do ponto da FEA aí em cima). Isso já vai ser explicado melhor. Fiquem
atentos para não se perderem, hein?

Em 2012 foram implantadas as linhas 8012/10 (circular 1) e 8022/10 (circular 2).
Essas linhas funcionam dentro da USP, e nós alunos não pagamos, pois elas
aceitam o bilhete USP (BUSP -- Retire logo o seu!!).

%REFTIME
Em 2021 houve uma grande mudança no trajeto dos circulares, porque agora eles
têm um ponto final no Portão 3 (e tecnicamente não são mais circulares, será que
o nome vai ficar?). Como essa mudança ocorreu durante a pandemia, até mesmo
veteranes podem ser surpreendidos, então prestem atenção nos mapas!

Para acompanharem as rotas dos circulares em tempo real, bem como ver pontos
importantes da faculdade, podem usar o aplicativo ``Guia USP", disponível para
Android e iOS.

Para os bixes que não gostam de ler (preferem imagens) colocamos o
mapa das três linhas nas próximas páginas.

\mapa{8012_ida.png}
\mapa{8012_volta.png}
\mapa{8022_ida.png}
\mapa{8022_volta.png}

\end{subsubsecao}

\begin{subsubsecao}{Como ir e vir do IME pelo metrô Butantã}

Como sabemos que vocês, bixes, chegam muito perdidos, e muitos pularam todas
essas informações sobre os pontos de ônibus e por isso podem não saber como
fazer seu trajeto, resolvemos colocar aqui um dos trajetos mais comuns que boa
parte de vocês vão usar.

Para chegar no IME a partir do metrô Butantã, vocês devem pegar o circular 1
(8012-10). Se vocês pegarem o circular 2 (8022-10), em algum momento vocês 
vão chegar no IME, mas o circular 1 é muito mais rápido.

O circular 1 vai fazer o seguinte trajeto: Faculdade de Educação, CEPEUSP,
Praça do Relógio Solar, Biblioteca Brasiliana e Faculdade de Economia,
Administração e Contabilidade (FEA). A partir daí, deverão descer no ponto da
FEA.

Quando estiverem indo embora para o metrô, vocês devem pegar o circular 2
(8022-10) no ponto da Oceanografia. CUIDADO: não peguem o circular 1 aqui
se estiverem querendo ir para o metrô Butantã!!! Lembrem que ele passa na
FEA logo no começo do trajeto (e provavelmente foi nesse ponto que vocês
desceram do ônibus quando chegaram). Se vocês pegarem o circular 1 (8012-10)
no ponto FAU I, também vão (em algum momento) chegar no metrô Butantã, mas
boatos dizem que os outros circulares vão mais rápido. Aí, basta ficar no
ônibus sentado (se conseguir um lugar!) até chegar no metrô Butantã, onde 
todos do ônibus vão descer.

\end{subsubsecao}

\begin{subsubsecao}{Como ir e vir do IME pelo P3}

{\bf Circulares durante dias de semana}

{\bf Circulares durante finais de semana}

Teste

\end{subsubsecao}

\begin{subsubsecao}{Como ir e vir do IME pela Cidade Universitária}

{\bf Circulares durante dias de semana}

{\bf Circulares durante finais de semana}

Teste

\end{subsubsecao}

\begin{subsubsecao}{Linhas intermunicipais}

Agora, se vocês moram mais longe ainda (outra cidade, outro estado, outro país,
outra dimensão...) e não querem ou não podem se mudar para São Paulo, existem
algumas linhas de ônibus fretados para cidades mais próximas (ou não). Se a
cidade de vocês não estiver aí, procurem se informar a respeito, pois não
significa necessariamente que não haja ônibus da USP para lá. Aí estão elas:

\begin{itemize}
  \item {\bf Empresa Urubupungá}\\
    Tel: 0800 11 4777
    Site: {\tt www.urubupunga.com.br}\\
    280BI1- São Bernardo do Campo (Centro)\\
    Cor: Cinza\\
    Onde pegar para sair da USP: ponto FAU II\\
    Para mais informações sobre a rota e a tarifa dessa linha, acesse
    \url{http://itinerario.urubupunga.com.br:8080/itinerario/ItinerarioLinha.aspx?lin=46\&emp=1}

  \item {\bf Fretados Jundiaí - USP}\\
    Viação MIMO\\
    Tel: (11) 4606-8222\\
    {\tt www.viacaomimo.com.br}\\
    Principais horários:\\
    Ida: 6h00; 6h30; 7h00; 12h50; 18h00 (na Rodoviária de Jundiaí)\\
    Volta: 11h40; 15h00; 16h40; 17h05; 18h10; 21h00; 22h49 (No ponto FEA)

  \item {\bf ABC}\\
    Osmar: (11) 94710-8604\\
    Marcelo Antonio: (11) 94719-1783

  \item {\bf Santos}\\
    Arca Turismo\\
    (11) 5928-7961\\
    \url{http://www.arcaturismo.com.br}

\end{itemize}

\end{subsubsecao}

\end{subsecao}


\begin{subsecao}{Horário dos Portões: Veículos no Campus}

Saibam por onde entrar na USP. Lembrem-se de terem sempre a carteirinha USP, ou
o e-Card no celular, ou o comprovante de matrícula (nos primeiros dias), além de
RG, pois esses documentos podem ser solicitados principalmente nos horários de 
entrada controlada.
\begin{itemize}
  \item {\bf Portaria 1 (P1):} R. Afrânio Peixoto. Funciona 24h por dia todos os
  	dias, mas a entrada é controlada para pedestres e carros nos seguintes
	horários: todos os dias das 20h às 5h, aos sábados após as 14h, aos domingos
	e feriados o dia inteiro. É por onde entram os ônibus municipais.

  \item {\bf Portaria 2 (P2):} Av. Escola Politécnica. De segunda à sexta,
  	acesso livre das 5h às 20h e controlado das 20h às 0h. Aos sábados, acesso
	somente para pedestres: liberado das 5h às 14h e controlado das 14h às 0h.
	Fechada aos sábados (para veículos), domingos e feriados. Única entrada para
	caminhões.

  \item {\bf Portaria 3 (P3):} Av. Corifeu de Azevedo Marques. Tem o mesmo
  	horário de funcionamento do P1 para os pedestres. Para os carros, tem o
	mesmo horário do P1 exceto aos domingos, feriados e madrugadas (0h às 5h),
	nas quais ele fecha, ao contrário do P1.

  \item {\bf Portaria P' (Plinha):} R. Eng. Teixeira Soares. Funciona de segunda
  	à sexta das 6h às 18h. Fechado de sábados, domingos e feriados.

  \item {\bf Portarias de pedestre (Mercadinho, São Remo, HU, CPTM e
  	Vila Indiana):} Funcionam de segunda à sexta das 5h às 20h, e de sábado
	das 5h às 14h. Acesso controlado de segunda à sexta das 20h às 0h.
	Nas portarias do Mercadinho e Vila Indiana, pode entrar em quaisquer
	outros horários, incluindo domingos e feriados, mas o acesso é controlado.

\end{itemize}

As linhas de ônibus municipais que passam por dentro da USP não operam aos
finais de semana, com exceção dos circulares, que entram na Universidade a
qualquer hora, e da linha 702U (Term. Pq. D. Pedro II - Cid. Universitária), que
funciona todos os dias das 5h às 0h (aproximadamente!). Se vocês vêm de carro,
saibam que a universidade dispõe de bolsões de estacionamento gratuito em torno
das Unidades.

\end{subsecao}
\pagebreak
\begin{subsecao}{Pontos de táxi}
Existem alguns pontos de táxi espalhados pela Cidade Universitária. As frotas
operam de segunda à sexta-feira, das 7h às 23h, além de aos sábados, das 7h às
17h. Eis suas localidades:

\begin{itemize}
\item Praça das Agências Bancárias\\
Fone: (11) 3091-4488

\item Praça da Reitoria\\
Fone: (11) 3091-3556

\item Hospital Universitário\\
Fone: (11) 3091-3536
\end{itemize}
\end{subsecao}

\end{secao}


\pagebreak

% Guia de jogos da Vivência ----------------------------------------------------
\input{guia_de_jogos.tex}

%FIXME GAMBIARRA para não quebrar página num lugar zuado
\pagebreak

% Dicas ------------------------------------------------------------------------
\input{dicas.tex}

%FIXME GAMBIARRA para não quebrar página num lugar zuado

% Melodias para a bixarada -----------------------------------------------------
\pagebreak
\begin{secao}{Músicas para a Bixarada}

\begin{subsecao}{Hino Cabeção}
\begin{verse}

Ó IME USP, meu amor\\
Para sempre campeão\\
Cálculo só me traz dor\\
E a REC emoção

Com a cachaça na mão\\
Para teus times vou torcer\\
Jogar com raça e coração\\
Vermelho e branco até morrer

ôôô ôôÔÔ ôôÔÔ ôôÔÔ

Eu vim aqui\\
Só pra gritar\\
Dá-lhe! IME! USP!
\end{verse}
\end{subsecao}

\begin{subsecao}{Musiquinhas de inter}
\begin{verse}

``O IME é o IME do IME"

``IME chupa IME lambe"

``IME ao contrário é M. M de Mackenzie"

``O IME é uma droga, vamos fumar o IME"

``Eu sou do IME, sou mais meu time\\ 
\,Não uso drogas e assisto anime"

\end{verse}
\end{subsecao}
\pagebreak

\begin{subsecao}{Músicas do BIFE}
\begin{verse}

(No ritmo de "Pelados em Santos")

''IME!''
''Estaremos contigo''
''Tu és minha paixão''
''Não importa o que digam,''
''sempre estarei contigo!''

''Minha camisa vermelha''
''E a cachaça na mão,''
''o BIFE me espera,''
''Vai começar a Festa!''

''Sha laiá laiá''
''Sha laiá laiá''
''IME do meu coração!''

\end{verse}
\end{subsecao}
\pagebreak

%\begin{comment}\begin{subsecao}{Caboclo da MAT}

% {\em Cantar como ``Faroeste Caboclo'' da Legião Urbana. Essa música é em memória 
% à FUVEST, quando Matemática Aplicada e POLI pertenciam à mesma carreira}

% \begin{verse}
% \footnotesize{
% Cheio de medo em setembro Joãozinho viu que seus dedos tremiam pra fazer a
% inscrição\\
% Deixou pra trás a namorada, a motoca, o futebol e as festinhas pra rachar na
% revisão\\
% Quando criança só pensava em ser engenheiro ainda mais com o dinheiro que
% sonhava em ter na mão\\
% Era o CDF lá do colégio onde estudava e todo mundo admirava o boletim desse
% cuzão\\
% Ia pra igreja só pra rezar pro seu santo pra pedir a sua ajuda pra prestar
% vestibular\\
% Sabia mesmo que ia ser barra pesada porque tinha muito japa pra tomar o seu
% lugar\\
% O ano todo se propôs a estudar, passava o dia sem ligar a televisão\\
% Nos feriados não ia viajar, ficava em casa treinando redação\\
% Fazia todos os exercícios da apostila e no fim de cada aula ia falar com o
% professor\\
% Às quinze horas ia pro laboratório ver as mitocôndrias da aula anterior\\
% Não entendia como o militarismo dominou nosso país por vinte anos de terror\\
% Ficou cansado de tentar achar resposta e desceu pra lanchonete pra afogar a sua
% dor\\
% E lá chegando foi tomar um cafezinho e encontrou um concorrente com quem foi
% falar\\
% E o concorrente aumentou seu desespero pois manjava muita coisa que ele tinha
% que estudar\\
% Dizia ele, eu vou prestar o ITA... Nesse país prova pior não há\\
% E se não der eu vou pegar engenharia, lá na POLI eu vou tomar o seu lugar\\
% E João não gostou dessa proposta, ele disse ``ai que bosta, eu tô passando mal''\\
% Ele ficou bestificado com a ideia de pegar lista de espera só depois do carnaval\\
% Meu Deus, é pior ainda, no ano novo eu posso estar lá na Mauá\\
% É brincadeira querer ser engenheiro e só descolar emprego em Taguatinga\\
% Na sexta-feira ele morria de vontade de correr pro banheiro se borrando de pavor\\
% E conhecia muita gente arrogante que passava do seu lado se dizendo um terror\\
% Ele estudava o relevo da Bolívia, função quadrática e modular\\
% E nos domingos então ele fazia tarefa mínima e complementar\\
% E Joãozinho até a morte se esforçava e o tempo mal sobrava pr'ele se alimentar\\
% E via às duas horas o Vestibulando que passava todas as dicas sobre o vestibular\\
% Mas ele não queria mais conversa e decidiu que em novembro era hora de rachar\\
% Ele pirou que precisava estudar tanto, virou um bitolado e começou a delirar\\
% E logo, logo os malucos da sua idade viram a calamidade, tem babaca novo aí\\
% E o nosso Joãozinho ficou louco e bateu em todos os japoneses dali\\
% Seus amigos preocupados com a sua sorte deram uma fita de rock pr'ele relaxar\\
% Mas de repente sob uma má influência dos boyzinhos lá do fundo começou a zoar\\
% Já na primeira fase ele penou e só passou porque o corte foi sessenta e três\\
% A demência tomou a sua mente: ``Vocês vão ver, eu vou pegar vocês!!!''\\
% Agora Joãzinho era fodido e estava decidido que não ia se dar mal\\
% Sacava toda a trigonometria e manjava de limites, derivada e integral\\
% Foi quando conheceu uma menina e de toda aquela zona ele se arrependeu\\
% Maria Lúcia era uma bitola linda e o coração dele pra ela o Joãzinho prometeu\\
% Ele dizia que devia estudar, pois engenheiro ele queria ser\\
% Maria Lúcia, pra sempre vou te amar, Engenharia com você quero fazer\\
% O tempo passa e um dia chega a hora de fazer segunda fase coitadinho do João\\
% E ele faz uma prova perigosa diz que espera uma resposta, pode ser um sim ou não\\
% Não vou correndo pra banca de jornal nem pra pátio do cursinho isso eu não faço
% não\\
% Pois eu prefiro ficar na minha casa esperando o resultado com o cu na mão\\
% Maria Lúcia vai comprar o tal jornal e logo após achar seu nome ela procura o de
% João\\
% Mas ela volta com tristeza no olhar, olha pra ele e diz ``você pegou a quarta
% opção''\\
% Você passou na sua quarta opção, você passou na sua quarta opção\\
% Bacharelado em Matemática é um tesão, eu vou sofrer as consequências como um cão\\
% Não é que Joãozinho estava certo, seu futuro era incerto mas foi se matricular\\
% Matriculou-se e no meio da zoeira descobriu que tinha muitos como ele no lugar\\
% Fez inscrição pro remanejamento e talvez no fim do ano transferência ia tentar\\
% E Joãozinho mantinha a esperança de um dia ir pra Poli estudar química\\
% Mas acontece que um tal Professorzzini terrorista de renome apareceu por lá\\
% Ficou sabendo dos planos de Joãozinho e decidiu que com suas notas ele ia se
% ferrar\\
% E ele teve que largar Cálculo 2 mesmo sabendo derivar e integrar\\
% E decidiu deixar Estat pra depois que o Moretin voltasse a lecionar\\
% Professorzzini, professor mais sem vergonha com sua prova enfadonha fez todo
% mundo dançar\\
% Desvirginava bixetes inocentes e o nabo era tão quente que nem dava pra sentar\\
% E Joãozinho há muito não via sua amada, e a saudade começou a apertar\\
% Eu vou pra Poli eu vou ver Maria Lúcia, já está em tempo de a gente se encontrar\\
% Chegando à Poli então ele chorou quando viu Maria Lúcia namorando um japonês\\
% Oh, Maria Lúcia, quanto que você mudou, que estrago que a Poli te fez\\
% Joãozinho era só ódio por dentro e então o japonês para um duelo ele chamou\\
% Amanhã às duas horas no Biênio, ou na Praça do Relógio, seja lá onde for\\
% E você pode escolher as suas armas: derivadas ou matrizes de qualquer versor\\
% Que eu provo que o sub-espaço nulo é o coração dessa piranha a quem jurei o meu
% amor\\
% E Joãozinho não sabia o que fazer quando escutou um papo lá no bandejão\\
% Onde falavam dum duelo que iam ver dizendo a hora, o local e a razão\\
% No sábado então às duas horas toda a Poli sem demora foi lá só pra assistir\\
% Um japa que botava pelas costas, encoxou Maria Lúcia e começou a sorrir\\
% Sentindo um ódio na garganta João olhou pros cabacinhos e pros trouxas a
% aplaudir\\
% E olhou pros pipoqueiros e as bancas de cachorro-quente que passavam por ali\\
% E se lembrou de quando era uma criança e de tudo que vivera até ali\\
% E decidiu entrar de vez naquela dança, se a Poli é um circo, e daí\\
% E nisso o céu abriu seus olhos e então Maria Lúcia ele reconheceu\\
% Ela queria fazer Álgebra 2 pra provar que a Poli não a emburreceu\\
% Politécnico, eu sou homem coisa que você não é, e não me contento em por nas
% costas não\\
% Some daqui filha da puta sem vergonha vai pra casa tocar bronha o seu destino é
% ser bundão\\
% E Joãozinho deu as costas para os dois, foi pra Pura onde encontrou o seu valor\\
% Maria Lúcia se arrependeu depois prestou Fuvest mas no IME não entrou\\
% E a todos declarava que o nosso Joãozinho era gênio que escapou de se foder\\
% Que na alta burguesia lá da Poli todo mundo é bunda mole ninguém sabe o que
% fazer\\
% E foi dar monitoria no cursinho pra avisar aos molequinhos pra não esquecer\\
% Ele queria era avisar toda essa gente engenharia é pra demente que só quer
% sofrer.\\
% }
% \end{verse} \end{comment} 
%\end{subsecao}

\end{secao}


% Utilidades -------------------------------------------------------------------
\begin{secao}{Utilidades}

\begin{subsecao}{Na WEB}

\url{https://usp.br} - Página da USP. Aqui vocês encontrarão notícias e eventos da
universidade, bem como informações gerais.

\url{https://ime.usp.br} - Página do IME. Nesse link, vocês poderão ver detalhes sobre
os cursos e obter informações sobre a faculdade.

\url{https://uspdigital.usp.br/jupiterweb} - Sistema JúpiterWeb. Aqui vocês vão
encontrar a grade horária e, mais tarde, poderão fazer matrícula nas matérias
que irão cursar no semestre. Além disso, podem acompanhar o pagamento de
bolsas e auxílios por ele e também ter acesso ao seu histórico escolar com
as suas notas.

\url{https://portalservicos.usp.br} ou \url{https://olimpo.usp.br} - Portal com todos os
serviços digitais da USP. Vocês vão reparar que a aba ``Graduação'' na verdade
é o JúpiterWeb disfarçado, então na prática vocês podem só acessar o link acima
que dá na mesma (só é um pouco menos bonito porque é um site mais antigo).

\url{https://edisciplinas.usp.br} - É bixes\dots Acham que vai ser essa moleza pra
sempre? Se vocês acham, estão muito enganados! Daqui a pouco vocês vão receber
uma senha para poder enviar seus EPs (vide glossário) nesse endereço\dots
(e não adianta fazerem chantagem emocional que o e-Disciplinas só vai aceitar
até 23h59\dots Não entendeu? Vocês vão entender\dots)

\url{https://instagram.com/spottedimeusp/} - Viu alguém interessante? Quer mandar uma
cantada pro crush, mas tem vergonha? Manda um spotted! Afinal, só o $x$ deve
ficar isolado. Ou achou alguém bacana e com gostos parecidos, que tal mandar um
spotted para fazer novas amizades?

\url{https://www.instagram.com/usp_spotted/} - Spotted da USP. Serve para a mesma
coisa que o Spotted do IME, só que é para USP toda.

% \url{ www.xkcd.com} - Webcomic sobre matemática e computação. Origem de muitas
% piadas que vocês ouvirão por aí.

E por fim\dots

\url{https://google.com.br} - Tudo o que vocês precisam tá no google. Se não
estiver, então não existe. Aproveitem e procurem o Código de Ética da USP (ou
baixem ele em \url{http://www.prg.usp.br/wp-content/uploads/CodigoEtica.pdf})

\end{subsecao}

\begin{subsecao}{Apps}
	
Alguns dos apps disponíveis na AppStore e PlayStore oferecidos pela própria USP
que podem ajudar. Para ver a lista completa, procure pelo desenvolvedor
``Universidade de São Paulo'' em qualquer uma das duas lojas.

{\bf e-Card USP} - Esse é o app mais importante nesse começo de ano, enquanto
sua carteirinha não chega (ver seção Cartão USP no Júpiter). O aplicativo é uma
carteirinha virtual que te dá acesso a tudo que o cartão físico permite,
inclusive comer no bandejão.

{\bf Cardápio+ USP} - Contém o cardápio dos bandejões e mais outros serviços
da SAS, como creche, saúde mental e moradia. Lá vocês também poderão conferir 
se o bandejão está fechado ou funcionando com o horário reduzido.

{\bf Guia USP} - Mapa dos principais lugares da USP, incluindo os institutos,
restaurantes, bancos e mais. Mostra a localização dos circulares em tempo quase
real, mas não é muito preciso e às vezes trava.

{\bf Campus USP} - Permite acessar a central da Guarda Universitária
diretamente no caso de uma emergência. Contém um mapa com as ocorrências
recentes de segurança.

{\bf Disque Trote USP} - Use esse app para denunciar caso sofra um trote
violento. Se preferir conversar mesmo, também pode ligar para o número de
telefone que está mais embaixo.

\end{subsecao}

\begin{subsecao}{Telefones}

{\bf Disque Trote:} {\tt 0800-012-1090} --
\underline{Não deixe de ligar se você sofrer algum trote violento!!}

{\bf SAS:} Seção de Passe Escolar: {\tt (11) 2648-1863} ou {\tt (11) 2648-1865}

{\bf Serviço de Graduação do IME (mais conhecido como Seção de Alunos):} {\tt (11) 3091-6149}

{\bf Superintendência de Tecnologia da Informação (STI):} {\tt (11) 3091-6400}

%{\bf Plantão de Cálculo e Álgebra Linear (serviço gratuito):} {\tt 3037-1773}
%Diziam que era o orelhão do ime, vai que isso volta, resolvi só comentar^^

\end{subsecao}
\end{secao}


\pagebreak
% Glossário --------------------------------------------------------------------
\begin{secao}{Glossário}

Esta parte é, sem dúvida, uma das mais importantes do guia. Aqui vocês
encontrarão todas as explicações para as maiores dúvidas do universo, apesar de
já sabermos que a resposta para a vida, o universo e tudo mais é 42. Com
certeza, após lerem este trecho sua vida vai mudar: vocês saberão, por exemplo,
porque o céu é azul e com quantos paus se faz uma canoa.

\begin{subsecao}{Sobre as matérias}

{\bf Teorema:} um teorema é uma afirmação que pode ser provada. Provar
teoremas é a principal atividade de quem estuda matemática. Deles surgem Lemas,
Corolários, Proposições, e tantas outras coisas que você só vai entender
completamente o que significam quando precisar escrever sobre eles, o que vai
acontecer logo, logo!

{\bf IC (Iniciação Científica):} grupo de estudantes que se dedicam a um
determinado assunto, sob orientação de alguém da área, paralelamente ao curso.
No caso de estudantes que queiram bolsas de estudo, é adotada uma abordagem
mais rigorosa.

{\bf P1, P2, P3:} \sout{provas} um dos maiores sofrimentos da graduação. 
A intensidade depende de quem ministra a disciplina.

{\bf SUB:} prova que você faz quando vai mal nas P1, P2 e P3, ou quando você
simplesmente não vai. Sua aplicação e utilidade depende de quem é responsável
pela matéria.

{\bf REC:} prova que você faz quando vai mal na SUB. Para fazer a REC, precisa
ter pelo menos média 3 nas outras avaliações.

{\bf DP:} matéria que você faz quando vai mal na REC.
\end{subsecao}

\begin{subsecao}{Sobre programas}

{\bf Computador:} objeto com vontade própria, sensível, que requer muito
carinho e atenção.

{\bf EP:} Exercício-Programa. Algo que você vai ter que fazer muitas vezes, e
vai dar muito trabalho.

{\bf GCC:} compilador mais recomendado para seus EP's, por suas inúmeras
qualidades. Atenção: ele ainda fará você se sentir incompetente.

{\bf Hello World:} um clássico da programação universal.

{\bf Segmentation Fault:} efeito computacional aleatório causado pela ``véspera
de entrega de EP''. Desenvolvido por Murphy.

{\bf Stack Overflow:} mensagem que aparece na tela do computador
quando ele se recusa a funcionar. Também é um site com todas as respostas
para todas as perguntas sobre programação: \url{https://stackoverflow.com/}.

{\bf Teorema Fundamental do EP:} ``O EP só funcionará no dia da entrega.'' Não
confunda com o Corolário 42 da Lei de Murphy: ``O EP só {\bf não} funcionará no
dia da entrega!''

{\bf Linux:} sistema operacional criado totalmente em linguagem C, graças a um
esforço mundial de milhares de experts em programação e informática, composto
por aproximadamente 7 mil arquivos e 5 milhões de linhas, e com o qual você não
tem capacidade para trabalhar.

{\bf Windows:} Vírus. Porém tão bem mascarado que parece até a coisa correta a
se usar.
\end{subsecao}

\begin{subsecao}{Sobre a USP}

{\bf CEPE:} Centro de Práticas Esportivas - lugar onde você poderá praticar
todos os esportes que quiser.

{\bf SAS:} Superintendência de Assistência Social. É onde estudantes fazem as
carteirinhas de passe de ônibus, EMTU e metrô, além de solicitar os auxílios
eteceteras que estão explicados na respectiva seção.

{\bf CRUSP:} Conjunto Residencial da USP. Se você se inscrever, torça para
pegar um apartamento num dos blocos já reformados, ou então torça para não
conseguir nenhum.

% {\bf Colmeia:} conjunto de favos.

% {\bf Favos:} um monte de prédios hexagonais encravados no meio do CRUSP.

{\bf CINUSP:} Cinema da USP, localizado no favo 4 da Colmeia. Toda semana ele
passa um filme de qualidade. Informe-se sobre a programação em
\url{https://www.usp.br/cinusp}.

{\bf Pelletron:} prédio da Física que na realidade é um acelerador de
partículas.

{\bf P1, P2, P3:} não confundir com as do item das matérias. Aquelas eram

{\bf as} P1, P2, P3; estes aqui são {\bf os} P1, P2, P3. São os portões da USP.
O P1, que é o principal, se localiza no cruzamento da Av. Afrânio Peixoto com a
Rua Alvarenga, além de ser por ele que entram os circulares vindos do metrô.
O P2 fica na Av. Escola Politécnica. Já o P3 fica na Av. Corifeu de Azevedo
Marques. Fique atento quanto aos horários dos portões, pois eles não ficam
abertos o tempo todo.

{\bf H.U. (Hospital Universitário):} para onde são mandades bixes que se
machucam na Semana de Recepção. Lá são realizados os tratamentos e experiências
com exposição à radiação, exposição a aspirantes a médicos, teste de
paciência/resistência a dor assim que pega a senha, alto nível de gesso no
estômago, queda espontânea (ou não) de cabelo etc.

{\bf Vet (Veterinária):} para onde são mandades biches que se machucaram na
Semana de Recepção.

{\bf Psico:} para onde são mandades bixes (e veteranes também) que querem
plantões psicológicos, porque saúde mental é importante e deve ser tratada com
seriedade. Ponto.

\end{subsecao}
\end{secao}


\pagebreak
% Considerações Finais ---------------------------------------------------------
\input{consideracoes_finais.tex}

%FIXME
%REFTIME
% O guia precisa ter um número de páginas totais múltiplo de 4, e a última
% página tem que estar vazia
\paginaembranco
\paginaembranco
\paginaembranco
\paginaembranco

\end{document}
